% Options for packages loaded elsewhere
\PassOptionsToPackage{unicode}{hyperref}
\PassOptionsToPackage{hyphens}{url}
\PassOptionsToPackage{dvipsnames,svgnames,x11names}{xcolor}
\PassOptionsToPackage{space}{xeCJK}
%
\documentclass[
  11pt,
  oneside]{book}
\usepackage{amsmath,amssymb}
\usepackage{iftex}
\ifPDFTeX
  \usepackage[T1]{fontenc}
  \usepackage[utf8]{inputenc}
  \usepackage{textcomp} % provide euro and other symbols
\else % if luatex or xetex
  \usepackage{unicode-math} % this also loads fontspec
  \defaultfontfeatures{Scale=MatchLowercase}
  \defaultfontfeatures[\rmfamily]{Ligatures=TeX,Scale=1}
\fi
\usepackage{lmodern}
\ifPDFTeX\else
  % xetex/luatex font selection
  \setmainfont[BoldFont=DejaVuSerif-Bold.ttf,ItalicFont=DejaVuSerif-Italic.ttf,BoldItalicFont=DejaVuSerif-BoldItalic.ttf]{DejaVuSerif.ttf}
  \setmonofont[BoldFont=DejaVuSansMono-Bold.ttf,ItalicFont=DejaVuSansMono-Oblique.ttf,BoldItalicFont=DejaVuSansMono-BoldOblique.ttf]{DejaVuSansMono.ttf}
  \ifXeTeX
    \usepackage{xeCJK}
    \setCJKmainfont[BoldFont=NotoSerifCJKsc-Bold.otf,AutoFakeSlant=0.15]{NotoSerifCJKsc-Regular.otf}
          \fi
  \ifLuaTeX
    \usepackage[]{luatexja-fontspec}
    \setmainjfont[BoldFont=NotoSerifCJKsc-Bold.otf,AutoFakeSlant=0.15]{NotoSerifCJKsc-Regular.otf}
  \fi
\fi
% Use upquote if available, for straight quotes in verbatim environments
\IfFileExists{upquote.sty}{\usepackage{upquote}}{}
\IfFileExists{microtype.sty}{% use microtype if available
  \usepackage[]{microtype}
  \UseMicrotypeSet[protrusion]{basicmath} % disable protrusion for tt fonts
}{}
\makeatletter
\@ifundefined{KOMAClassName}{% if non-KOMA class
  \IfFileExists{parskip.sty}{%
    \usepackage{parskip}
  }{% else
    \setlength{\parindent}{0pt}
    \setlength{\parskip}{6pt plus 2pt minus 1pt}}
}{% if KOMA class
  \KOMAoptions{parskip=half}}
\makeatother
\usepackage{xcolor}
\usepackage[paperwidth=176mm,paperheight=250mm,left=23mm,right=23mm,top=24mm,bottom=24mm]{geometry}
\usepackage{longtable,booktabs,array}
\usepackage{calc} % for calculating minipage widths
% Correct order of tables after \paragraph or \subparagraph
\usepackage{etoolbox}
\makeatletter
\patchcmd\longtable{\par}{\if@noskipsec\mbox{}\fi\par}{}{}
\makeatother
% Allow footnotes in longtable head/foot
\IfFileExists{footnotehyper.sty}{\usepackage{footnotehyper}}{\usepackage{footnote}}
\makesavenoteenv{longtable}
\setlength{\emergencystretch}{3em} % prevent overfull lines
\providecommand{\tightlist}{%
  \setlength{\itemsep}{0pt}\setlength{\parskip}{0pt}}
\setcounter{secnumdepth}{-\maxdimen} % remove section numbering
\ifLuaTeX
\usepackage[bidi=basic]{babel}
\else
\usepackage[bidi=default]{babel}
\fi
% get rid of language-specific shorthands (see #6817):
\let\LanguageShortHands\languageshorthands
\def\languageshorthands#1{}
% Shared print conventions; the two books keep separate prose and pagination.
\usepackage{fvextra}
\usepackage{xurl}
\usepackage{etoolbox}
\usepackage{fancyhdr}
\setCJKmonofont{NotoSansMonoCJKsc-Regular.otf}[AutoFakeBold=1.5]
\setlength{\emergencystretch}{3em}
\setlength{\parindent}{1.8em}
\setlength{\parskip}{0.32em}
\linespread{1.12}
\renewcommand{\contentsname}{目录}
\renewcommand{\figurename}{图}
\renewcommand{\tablename}{表}
% Reserve enough room for three-digit page numbers in the mathematics contents.
\makeatletter
\renewcommand{\@pnumwidth}{2.4em}
\renewcommand{\@tocrmarg}{3.4em}
\makeatother
\DefineVerbatimEnvironment{verbatim}{Verbatim}{breaklines,breakanywhere,fontsize=\small}
\AtBeginEnvironment{longtable}{\small\setlength{\tabcolsep}{3pt}\renewcommand{\arraystretch}{1.15}}
\pagestyle{fancy}
\fancyhf{}
\fancyfoot[C]{\thepage}
\renewcommand{\headrulewidth}{0pt}
\fancypagestyle{plain}{\fancyhf{}\fancyfoot[C]{\thepage}\renewcommand{\headrulewidth}{0pt}}
\Urlmuskip=0mu plus 1mu
\widowpenalty=10000
\clubpenalty=10000
\ifLuaTeX
  \usepackage{selnolig}  % disable illegal ligatures
\fi
\IfFileExists{bookmark.sty}{\usepackage{bookmark}}{\usepackage{hyperref}}
\IfFileExists{xurl.sty}{\usepackage{xurl}}{} % add URL line breaks if available
\urlstyle{same}
\hypersetup{
  pdftitle={客体、演化与意义},
  pdfauthor={原始思想与论文作者：Degui Qian; 本版：中文协作修订稿（ChatGPT 辅助分析与文字修订）},
  pdflang={zh-CN},
  colorlinks=true,
  linkcolor={MidnightBlue},
  filecolor={Maroon},
  citecolor={Blue},
  urlcolor={MidnightBlue},
  pdfcreator={LaTeX via pandoc}}

\title{客体、演化与意义}
\usepackage{etoolbox}
\makeatletter
\providecommand{\subtitle}[1]{% add subtitle to \maketitle
  \apptocmd{\@title}{\par {\large #1 \par}}{}{}
}
\makeatother
\subtitle{一种开放的统一自然哲学}
\author{原始思想与论文作者：Degui Qian \and 本版：中文协作修订稿（ChatGPT 辅助分析与文字修订）}
\date{2026-10-30}

\begin{document}
\frontmatter
\maketitle

{
\hypersetup{linkcolor=MidnightBlue}
\setcounter{tocdepth}{1}
\tableofcontents
}
\mainmatter
\textbf{基于 Degui Qian 原论文的中文协作修订稿}\\
2026 年 10 月 30 日

原始思想与论文作者：\textbf{Degui Qian}。本书以《A Universal Evolutionary Object Theory (UEOT): A Framework and Research Program for a Unified Science》为思想来源，在后续核心数学研究的约束下重新组织论证，并由 ChatGPT 协助完成分析与文字修订。原论文及其作者说明保留为历史来源。本书的概念论证独立展开，技术命题、证明及适用条件由配套数学著作承担。\footnote{本书是基于原始论文与后续研究的协作修订文本。原作者署名遵循原论文第 43 页的贡献说明；本版的新增篇章结构、规范论证及数学反馈按附录 B---D 定位。它不将原作者未逐项确认的新表述追溯归入原稿。}

\hypertarget{ux524dux8a00ux4eceux7edfux4e00ux7684ux613fux671bux8d70ux5411ux53efux4ee5ux5171ux540cux751fux6d3bux7684ux601dux60f3}{%
\chapter{前言：从统一的愿望走向可以共同生活的思想}\label{ux524dux8a00ux4eceux7edfux4e00ux7684ux613fux671bux8d70ux5411ux53efux4ee5ux5171ux540cux751fux6d3bux7684ux601dux60f3}}

一个人在河边看见流水，走进城市看见制度，回到家中感到牵挂，最后又独自追问生活的意义。他所经历的似乎是几个不同的世界：物质受规律支配，生命为生存而活动，社会通过规则协调，人则按照理由作出选择。我们已经学会分别描述这些世界，却常常不知道怎样把它们放回同一个世界。物理解释越成功，目的与价值是否就越显得多余？意义越贴近个人经验，它是否就越难拥有公共的真实性？

原始 UEOT 论文以极大的雄心回应这个分裂。它试图用客体的组织、演化中的双重关切以及多层级相互作用，把物质、生命、心灵和意义联系起来。它的价值首先在于坚持一个问题：如果人确实生活在自然之中，那么人的目的、承诺与自我理解，也应当能够在自然世界里找到位置。我们不应当一边承认所有活动具有现实条件，一边把最重要的人类经验逐出现实。\footnote{原著第 1---3、40---41 页。关于物质、生命、心灵与意义的共同世界，是本书承接的总体问题；原著的普适和终极表述在本版中按不同论证责任展开。}

本书继续这一追问。它主张，理解世界的一条有力道路，是考察相对稳定的组织怎样在变化中形成、怎样通过关系维持、怎样获得与环境打交道的能力，以及这种能力怎样发展为有理由的行动。但这条道路的说服力，不能仅来自重复``万物统一''。我们需要说明某种统一究竟消除了什么隔阂，保留了什么差别，并且使哪些原来难以提出的问题变得可以回答。

例如，身体更新了许多材料，人为什么仍认为昨天的承诺约束今天的自己？一群互相影响的参与者，在什么意义上形成了一个组织？能够预测某个系统，是否意味着已经知道它由什么构成？学会维持自身的生命，是否因此获得了值得追求的全部目的？一个社会长期稳定，是否足以说明它公正？这些问题彼此关联，却不能用同一个答案替换。统一思想的成熟，表现在它能够解释这些差别怎样相接。

本书提出的立场可以称为一种\textbf{关系性的客体实在论}。客体是真实世界中能够承受某些变化、保存某些差别并参与相互作用的组织。我们的提问方式决定了哪些差别值得追踪，但差别能否经受干预、预测和共同检查，并不由提问者任意决定。这个立场也具有\textbf{过程性的自然主义}：持续存在是变化的成就，心灵和价值活动有其身体、历史与社会条件。自然主义在这里是一项解释承诺：尽可能在共同世界中理解我们，而不是预先宣布所有解释都必须缩成同一种句法。

这样的哲学需要两种勇气。一种是提出大问题，承认局部知识之间仍有深刻的联系值得寻找。另一种是在论证要求时改变原来的表达。把``所有稳定存在必有几何闭环''改写为``应当辨明维持某种组织所需的因果依赖''，并没有放弃对存在的追问；它让这一追问更接近可以争论的对象。把``每个客体具有唯一终极目的''改写为``目的解释需要指出何种机制和评价关系赋予行动方向''，也没有把目的驱逐出自然；它使目的免于成为任何结果都能事后套用的称号。

以下各编沿着一条连续路线展开。第一编讨论变化、关系与客体；第二编讨论开放、修复、记忆和行动可能性；第三编讨论目的、价值与自我修订；第四编讨论组合、尺度与演化；第五编讨论心灵、自由、道德及社会；第六编回到方法、宇宙统一与真理。读者可以不阅读配套数学著作而理解这里的论证。需要核对形式化背景时，可使用书末的简短对应表。

本书所追求的统一，不要求所有声音归于一致。一个能够容纳真实分歧、解释自身修正并接受世界反驳的思想，比一个宣布再无问题的体系更接近共同求真的生活。

\part*{\phantomsection\addcontentsline{toc}{part}{第一编　变化中的客体}第一编　变化中的客体}

\hypertarget{ux7b2cux4e00ux7ae0-ux4e3aux4ec0ux4e48ux8fd8ux8981ux8ffdux95eeux7edfux4e00}{%
\chapter{第一章　为什么还要追问统一}\label{ux7b2cux4e00ux7ae0-ux4e3aux4ec0ux4e48ux8fd8ux8981ux8ffdux95eeux7edfux4e00}}

设想一座桥。工程师讨论它的承重，城市规划者讨论它连接的人群，居民讨论它怎样改变上学与就医的路程。三种描述并不互相排斥，却也没有一种可以单独完成另外两种的工作。桥的用途需要物理结构支撑，结构的修建又取决于制度与目的。若把所有描述都压成材料运动，我们会失去``为什么造这座桥''的问题；若只谈公共意义，我们又无法判断它是否会坍塌。

统一之所以值得追问，是因为现实后果穿过学科边界。一项理由促成一项制度，一项制度改变物质安排，物质安排重新塑造人的机会。解释若在边界处突然中断，我们便容易把自己尚未说明的部分当作另一个世界的神秘力量。统一思想首先要求相邻解释之间有可以检查的接合：一个层面的输入，如何成为另一个层面的条件；某种结论，在转换描述后还保留多少内容。

有人会反对：世界也许根本不统一，追求统一只是人的审美偏好。这个反对有其力量。简洁使理论易于掌握，却不能保证自然必须服从简洁。但是，统一的要求可以从实践中获得更温和的根据。我们在不同活动中与同一世界相遇，行为的后果不会因为换了术语而停止传播。把这些后果连接起来，是解释责任的一部分。由此可以要求相容性，却还不能要求所有规律来自一条终极公式。

因此，需要区分世界的共同性与理论的单一性。共同性是说，桥的物理过程和居民的行动没有生活在互不相干的宇宙里；单一性则是说，存在一套特定理论可以充分解释所有这些现象。前者是本书采取的现实主义立场，后者是一项需要逐步取得成果的研究企图。我们可以坚持前者而对后者保持开放。

另一种反对来自成功的分工：既然各门科学已经有效工作，何必寻找更一般的概念？回答在于，专业分工留下了一些不能由分工本身解决的问题。一个系统何时算作一个对象？何时是同一个对象发生变化，何时是另一个对象出现？控制、修复、选择和理解之间有什么共同结构？这些问题会在不同学科重复出现。若能够提出清楚的共同语言，就能减少概念偷换，也可能把一处已有的推理移用于另一处。

共同语言仍需接受局部现实的抵抗。一个能描述桥、身体和城市的概念，不能因为适用范围广，就取消材料、生命史和政治制度的区别。本书把统一理解为一种逐层完成的解释关系：有些概念共享，有些机制专属，有些桥接成立，有些桥接需要重写。其成败应当由解释带来的认识增量判断，而不是由纳入了多少宏大名词判断。\footnote{原著第 1---3 页；核心第 0、28、31 章。本章关于统一的哲学理由由正文论证，不把模块可组合性提升为所有自然理论已经统一。}

\hypertarget{ux7b2cux4e8cux7ae0-ux4e8bux7269ux600eux6837ux5b58ux5728ux4e8eux8fc7ux7a0bux4e4bux4e2d}{%
\chapter{第二章　事物怎样存在于过程之中}\label{ux7b2cux4e8cux7ae0-ux4e8bux7269ux600eux6837ux5b58ux5728ux4e8eux8fc7ux7a0bux4e4bux4e2d}}

日常语言把事物与变化分开：先有一个杯子，然后杯子变热、破损或被移动。这种说法方便，却容易使我们误以为事物有一个完全不依赖过程的内核。换一种眼光，杯子的形状需要材料结构持续承受环境作用；一个旋涡需要流动；一个组织需要人不断履行规则。它们的持续方式不同，但都不能脱离实际发生的过程来说明。

过程观点并不要求我们取消事物。恰恰相反，它要解释事物为什么能够成为稳定的谈论对象。若世界只有互不相关的瞬间，我们无法把一次观察与下一次观察连接起来。客体的出现，意味着某些差别可以跨越一段变化被追踪：改变位置后它仍能承担相近作用，替换部分材料后它仍保存某种组织，受到允许范围内的扰动后它仍具有可以预期的回应。

这里的``某些''十分重要。一块冰可以保持材料连续性而失去形状；一个机构可以保持名称而改变职能；一个人可以改变信念而继续承担过去的责任。客体性不是把所有属性都永久冻结。它要求说明，当前问题中哪些属性承担了追踪的任务，以及这种追踪怎样由实际过程支撑。

假如有人问：这种说法是否把客体降格为观察者的方便分类？可以用一个简单情形回答。两个人可能为了不同任务，把同一片森林划成不同区域；但任何一方都不能仅凭命名，使一条火带停止传播。分类服务于问题，传播关系却反过来限制分类的用处。对象描述具有选择性，并不意味着对象的全部现实性来自选择。

一个更小的模型也能说明这种客观性。设一个装置有表示运行条件的可见位和一个内部位；两种动作分别使下一步满足或不满足运行条件，而内部位重新取值。如果内部位不影响这些动作的后果，把四个完整状态按可见位分成两类，就能准确回答运行与控制的问题。若改问内部位的具体内容，原描述便不再充分。对象描述因问题而异，而它能否保留行动后果，仍受同一过程严格约束。第十四章将回到这个模型，考察这种准确描述还能怎样支持持续与目的的判断。

有限的描述也不规定整个观察世界必须有限。一个模型只在有限范围内运行，其他解释仍可能为范围外的情况保留概率。若为了方便而删除这些情况，再宣布模型与解释完全一致，所谓简化就改变了问题。合理的压缩可以保留一个明确的``范围之外''，使遗漏仍承担其概率后果；它未必继续区分外部情况彼此之间的全部差异。清楚说明保留了哪一种比较，才能使有限表示成为认识世界的工具。

过程观点因此同时反对两种极端。一种极端把事物看成不受关系影响的孤立实体，于是无法解释形成与消失。另一种极端把事物看成任意叙述，于是无法解释为什么有些分类可以预测和干预，另一些会失败。可行的中间立场，是把客体看作真实过程中的相对稳定组织：它的边界、时间尺度和重要属性可以随问题改变，但其成立需要世界提供足够的连续性与约束。

``稳定''也不等于静止。一个维持体温的身体不断交换物质，一个履行承诺的团队不断调整分工，一项持续有效的规则不断被具体执行。稳定可以是某些状态保持不变，也可以是变化方式保持可辨认。我们把这些不同情形放在一起，是为了理解持续存在的共同问题，而不是要求它们具有相同的微观机制。

由此，存在与生成不再是非此即彼。事物通过生成过程取得相对稳定，稳定组织又约束后来可能发生的生成。一个客体既是过去过程的结果，也是未来过程的条件。理解客体，要同时向后追问它怎样形成，向前追问它允许什么继续发生。\footnote{核心第 2---4、8、11---12 章，分别处理过程、表示、预测身份、持续与历史运输；数学著作第 36.8 节说明有限表示怎样在明确的一侧支撑条件下保留全部逃逸质量及相应概率比较。本文的``相对稳定组织''是综合这些区别的哲学用语。}

\hypertarget{ux7b2cux4e09ux7ae0-ux5173ux7cfbux8fb9ux754cux4e0eux5171ux540cux4e16ux754c}{%
\chapter{第三章　关系、边界与共同世界}\label{ux7b2cux4e09ux7ae0-ux5173ux7cfbux8fb9ux754cux4e0eux5171ux540cux4e16ux754c}}

一扇门既分开两个空间，又使两个空间以特定方式相连。客体的边界也常具有这种双重作用。身体的边界控制交换，制度的边界规定参与资格，仪器的边界决定哪些信号进入测量。若只把边界理解成外壳，就会遗漏这种选择与转换的功能。

关系观点强调，客体的许多能力来自它与其他事物的连接。钥匙的开锁能力不单属于钥匙的形状，也涉及锁的结构；一种语言表达的意义需要使用习惯与听者；某个组织的权力依赖其他参与者的承认与配合。不过，从这里不能推出``一切都是关系，所以关系两端无须说明''。钥匙仍须由具体材料实现，承认也须由实际的人和制度过程承担。关系改变我们对两端的理解，却不替我们省去对实现的解释。

划定边界时，可以提出三个相互补充的问题。第一，哪些变化可以从外部输入，而内部仍维持所讨论的组织？第二，哪些内部差别会改变对外响应？第三，哪些干预会使这种响应方式失效？这三个问题分别涉及开放、可辨识性与构成依赖。它们可能指向相近的边界，也可能指向不同的边界。一个数据库的访问权限边界，就不必与保存其数据的物理设备边界重合。

``边界依赖问题''容易被误读为``现实没有边界''。实际上，问题选择了检验方式，检验揭示了现实联系。若为了预测室内温度，把建筑和加热系统作为一个对象很有用；若为了维修管道，必须进一步区分子系统。两种描述可以同时真实，因为它们追踪不同而相容的依赖。相反，若某种划分总把决定结果的交换排除在外，它就需要修改。

我们也应区别尚未观察到的关系与根本不存在的关系。有限的感受能力并不规定世界的全部范围。一个客体可能通过中间过程受到影响，而没有直接识别影响的能力。对人的认识而言，这意味着共同世界大于当前经验；对科学而言，这意味着描述接口时应当允许遗漏被后来发现。

关系主义最有力的结论，并非所有东西都相互决定，而是解释必须给关系以准确的作用位置。有些关系传输信息，有些约束行动，有些维持材料，有些只是统计关联。只有区分这些作用，我们才能避免从相关性跳到因果性，从依赖跳到同一，从合作跳到整体利益。

客体的边界由此成为可讨论、可修正而又受现实约束的概念。它帮助我们组织共同经验，也提醒我们：任何一个视角能够清楚看到的世界，都可能把另一种重要联系留在边界之外。\footnote{核心第 6 章的开放、环境接口、因果方向与边界条件，第 13 章的访问与物理毁损区别。}

\hypertarget{ux7b2cux56dbux7ae0-ux3c9ux7ec4ux7ec7ux5982ux4f55ux7ef4ux6301ux81eaux5df1ux7684ux6761ux4ef6}{%
\chapter{第四章　Ω：组织如何维持自己的条件}\label{ux7b2cux56dbux7ae0-ux3c9ux7ec4ux7ec7ux5982ux4f55ux7ef4ux6301ux81eaux5df1ux7684ux6761ux4ef6}}

设想一个需要供电、控制和冷却的装置。供电使控制得以运行，控制调节冷却，冷却又维持设备能够继续工作。这里确有一种回返的依赖。理解它，不只是在纸上画出一个环，而是辨明哪些作用共同维持了继续活动的条件。这是本书保留 Ω 思想的核心理由。

原论文用闭环表达持续存在的内在组织。这种直觉抓住了被静态描述忽略的事实：许多系统的秩序必须被不断实现。其活动一方面消耗条件，另一方面重新建立条件。生命尤为鲜明，但类似问题也出现于技术系统和制度。\footnote{原著第 4---5 页关于 Ω 与存在的主张；原著第 37---38 页的拓扑数据提案。本章区分轨迹形状、因果回路与构成依赖。}

然而，组织闭合的含义需要从作用关系中获得。一个随外部节拍摆动的装置，轨迹可以形成漂亮的循环，却未必自行维持驱动条件；一个在稳定状态附近调节的系统，宏观轨迹未必呈现显著的几何圆圈，关键反馈仍可能真实存在。轨迹的形状、因果图中的回路、机制之间的再生依赖，回答的是不同的问题。

怎样把 Ω 变成有解释力的概念？可以从组织的失效方式入手。先说明要保持什么，再询问：移除哪些作用会破坏这种保持？如果两条独立通路都能供给必要资源，删去其中一条可能无碍，同时删去两条才会失效。组织的完整性因此不只是``有没有环''，还包括冗余、替代、瓶颈以及它们之间的组合。一个复杂系统可能有许多几何环，却只有一个脆弱的共同支撑；也可能外观简单，却具有多重恢复通路。

移除作用的实验同样需要解释。有时删除一个部件会解除抑制，使备份开始工作。此时，``删得越多必然越坏''并不成立。我们应当描述具体的响应，而不是要求世界迁就一个方便的单调图景。所谓完整性，始终要与所考察的组织能力、干预种类及时间窗口相连。

这种相对性不会使 Ω 失去本体论意义。它揭示的是：一个客体之所以能够持续，存在一组不能任意破坏的现实依赖。我们可以从不同任务辨认不同依赖，却不能使所有依赖都消失。客体的现实性，部分表现在它承受变化的能力和它不能承受变化的方式中。

对统一哲学而言，Ω 因而是一项积极的研究要求：寻找持续组织的构成依赖，并说明这些依赖如何被保存、替代和重建。它允许我们比较不同系统的维持方式，同时把具体机制留给相应的解释。一个符号可以引导多种研究；真正连接它们的，是共同问题及经得起检查的桥接。\footnote{核心第 12---14 章；配套数学著作第 32---34 章。极小失效表示需要相应有限、单调及干预解释条件；正文没有将这些条件规定为所有对象的普遍本质。}

\hypertarget{ux7b2cux4e94ux7ae0-ux540cux4e00ux4e2aux5ba2ux4f53ux8eabux4efdux662fux4e00ux6bb5ux600eux6837ux7684ux5386ux53f2}{%
\chapter{第五章　同一个客体：身份是一段怎样的历史}\label{ux7b2cux4e94ux7ae0-ux540cux4e00ux4e2aux5ba2ux4f53ux8eabux4efdux662fux4e00ux6bb5ux600eux6837ux7684ux5386ux53f2}}

一把经过多次维修的椅子，一家更换全部员工的机构，一个不断形成新记忆的人，都使``还是同一个吗''变得困难。如果身份只取决于材料相同，许多通常的持续关系会消失；如果只取决于功能相同，两个同时存在的复制品又会被认作同一个。身份需要比材料或功能单独提供的内容更多。

可以先区别相似、等价与历史同一。相似表示某些属性接近；等价表示在某类问题上无法区分；历史同一则要求一条有方向的延续关系。两份内容相同的文件在阅读任务中等价，但它们可以来自不同生成过程。一份文件经过授权修订，内容已不同，却仍属于同一份文档的历史。问题决定哪种关系重要，论证则必须明确使用哪一种。

历史身份的支撑不必是一条永远不变的属性。它可以是一连串有依据的传递：材料联系、组织保持、记录继承、责任承接以及双方认可的规则。每一步都要说明为什么后来的状态延续了所追踪的对象。连续传递有时能够容纳很大的总变化，但也可能在分裂、融合或替换中遇到真正的歧义。

设想一个研究团队分成两个团队，双方都保留原来的成员、资料和研究目标。若原来的名称只能归属一方，命名规则可以作出决定；但这个决定不会消除两条实际继承关系。此时更好的描述，是区分法律身份、研究谱系与组织延续，而不是期待一个抽象``同一''替我们解决所有问题。概念的精确可以表现为保留合理的多重关系。

这对人的身份也有启发。一个人能够改变信念而继续承担责任，是因为责任所依赖的历史关系通常比当前信念更宽。与此同时，严重的能力变化、强制处境或记忆损伤，可能改变合理归责的方式。承认这种复杂性，有助于我们既不把人冻成不变本质，也不因变化而抹去他与他人的历史联系。

修复特别清楚地展示了身份的方向性。修复工作的目标通常是恢复这一个客体的某种能力，并保存所要求的连续关系；制造一个功能相同的替代品则可能完成另一项任务。替代品可以同样有价值，但其价值不需要通过假装历史完全相同来获得。繁殖又是第三种关系：后代继承某些组织，同时作为另一个个体进入世界。

把身份理解为受约束的历史，使我们能够解释成长与坚持为何能够并存。值得维持的，不总是过去每一项状态，而可能是使承诺、能力与共同生活得以延续的关系。一个客体的过去不是一枚静止标签；它通过可以说明的传递，继续参与现在。\footnote{核心第 11、13、15 章；后续客体角色、形成与绑定、谱系身份的精确关系见数学著作第 34 章。}

\part*{\phantomsection\addcontentsline{toc}{part}{第二编　开放、记忆与行动的可能}第二编　开放、记忆与行动的可能}

\hypertarget{ux7b2cux516dux7ae0-ux5f00ux653eux7684ux5b58ux5728ux4e0eux6709ux9650ux7684ux8eabux4f53}{%
\chapter{第六章　开放的存在与有限的身体}\label{ux7b2cux516dux7ae0-ux5f00ux653eux7684ux5b58ux5728ux4e0eux6709ux9650ux7684ux8eabux4f53}}

一个需要呼吸、摄取、休息的生命，无法把自身保存为与世界隔绝的完整物。它必须让某些东西进入，让某些东西离开，并在这种交换中保持组织。开放因此不是对已经完成的存在所作的附加，而是许多存在方式本身的条件。

但开放也带来暴露。能接收养分的接口可能接收有害物，能获得消息的渠道可能传来误导，能与他人合作的关系也可能形成依赖。开放的意义不在于交换越多越好，而在于交换怎样参与组织的维持与发展。保护边界和扩大联系，都可能在不同处境中成为合理行动。

原文区分较为稳定的客体与高度依赖环境交换的客体，并用开放程度描述生命周期。这个区分可以引导比较：有些结构一旦形成，较少依赖持续调节；另一些组织必须不断工作，才能继续保持自身。然而，开放程度很难成为把一切事物排列成单一序列的尺子。交换什么、以什么速度、经过哪些接口、对何种能力不可替代，都可能改变比较的结果。\footnote{原著第 4---5、9---11 页。正文将开放作为需要按内容与接口说明的组织问题，而非已确认的单一生命周期规律。}

对生命而言，身体使这种开放具有一个不可随意移动的中心。身体既是行动的资源，也是损伤会发生的地方。我们可以放弃一个任务，不能同样轻易地放弃承受任务后果的身体。正因为行动有这种负担，世界才不是无差别的信息集合。某些变化意味着继续活动的可能，某些变化意味着能力被永久削弱。

由此可以理解一种基础的利害关系：对维持着某种组织的存在者而言，条件的改善与破坏具有不同后果。这种差别为生命的功能解释提供了现实根据。器官、调节和行为可以围绕它们怎样参与组织的延续来理解。但从这种基础利害关系走到完整的道德价值，还需要说明谁值得被关切、怎样处理冲突以及什么理由可以公开成立。身体使价值问题有了真实承受者，却不独自给出全部答案。

有限性也意味着必须区分平均上可维持与这一次能够承受。一个装置若平均获得足够资源，仍可能在某次资源到来之前已经停止工作。一个人长期来看拥有足够收入，也可能在急需照护的时刻得不到支持。若时间安排、储备与不可逆损伤被平均值遮蔽，关于``可持续''的判断就会失去重要内容。

身体的哲学意义，正在于把抽象可能性带回有时间、有地点、有代价的现实。行动不是在一张无限菜单上无损地挑选。我们通过行动消耗条件、改变环境，也改变后来还能做什么。开放生活的智慧，应当同时关心继续交流的能力和承受交流后果的能力。\footnote{核心第 8---10、22 章；数学著作第 34---35 章关于资源、有限路径、实际生存和可信修复的区别。均值与路径承受能力的例子用于解释这一概念差异。}

\hypertarget{ux7b2cux4e03ux7ae0-ux4feeux590dux751fux547dux4e0eux81eaux521bux751f}{%
\chapter{第七章　修复、生命与自创生}\label{ux7b2cux4e03ux7ae0-ux4feeux590dux751fux547dux4e0eux81eaux521bux751f}}

一个物体被外力修好，与一个组织能够参与自身修复，具有不同意义。后者把某些维持条件纳入了自己的活动：发现损伤、调用资源、重建通路、检查恢复是否有效。自我维持由此获得了更丰富的内容。客体不仅延续过去的运作，还对运作受到威胁作出有组织的回应。

然而，``自己修复自己''不是从无到有地创造一切。修复需要尚可工作的探测、执行和协调能力，也需要某种方式区别成功与失败。若修复程序完全失去解释它的机制，即使程序文本仍然存在，修复也无法仅凭文本发生。若所有恢复材料都已不可获得，正确步骤同样可能无法实现。这些支撑并非修复论的缺点，而是解释修复怎样发生的内容。

设想一台能够从备份恢复控制程序的设备。我们可以合理地说，它恢复了自身的一项能力。但是，读取备份的部件、执行恢复的规则以及程序对应现实动作的机制，都参与了这一成功。如果后来要说明这些支撑也怎样被形成，就需要进一步研究它们的来源。把形成问题与恢复问题分开，有助于看到生命组织的层次，而不是把所有困难隐藏在``自创生''这个名称里。

程序的来源问题还能从信息角度看得更清楚。两份局部状态记录可能完全相同，而程序的含义随着未被记录的解释约定或形成历史改变。只读取这份局部记录，就无法区分两种要求。要可靠恢复，必须保留足够的区别，或从其他通道取得它们。这说明恢复依赖信息与实现的组织，并不说明程序能够脱离物理条件存在。

自创生的哲学吸引力，在于它试图突破把生命理解成零件集合的方式。零件的作用需要置于它们共同维持的组织中，组织又通过具体过程产生和替换零件。这里有一种循环解释，但并非必然的逻辑循环。我们可以同时说明各个过程在不同时间怎样依赖，并通过扰动检验这种依赖。不可接受的循环，是把``能自我产生''既当作结论又当作唯一理由；可以解释的循环，是一个在时间中运行并有现实条件的再生网络。\footnote{原著第 4---5 页的自维持思想；核心第 9、26 章；数学著作第 34---35 章关于形成、程序来源及同源修复。程序来源的不足是信息与实现条件的问题。}

这项来源要求已有一个可以逐步检查的数学实现。有限寄存器保存一个目标程序；解释器读取实际寄存器，决定保持还是写入。环境故障会破坏目标位，程序本身按模型约定受到保护。修复、失败、费用与实际产生的子代保存在同一记录中，真实出生才推出子代继承程序并拥有下一次修复能力。这里的进展，是把``程序能够实行''和``孩子具有修复底座''从独立保证改成执行结果；可靠解释器、保护和继承的模型条件仍然公开。它没有使所有支撑都从自身产生。\footnote{数学著作第 39.1---39.3 节；RegisterRepairSource、RegisterFaultSource、RegisterBirthMaintenance 的实际执行及子代消费者。故障载体进入是修复能力，不是永远健康。}

有人可能担心：只要列出足够多的支撑，任何机器都能被称为生命。这个担心提示我们，生命概念不应由某一项孤立能力决定。自我维持、代谢式交换、发展的组织连续性、繁殖与遗传等能力可以形成相互关联的组合，而不同研究可能关注不同组合。把生命视为具有内在联系的能力群，比强迫一条简单边界承担全部判断更有解释力。

这也影响我们理解衰老与死亡。死亡不只是某个数值越过阈值，也可能是维持和恢复之间的关系无法再组织起来。某个局部部件仍然活动，并不足以说明整体生命已经延续；材料被其他生命利用，也不等于原来的个体继续存在。承认这种区分，使我们能够同时尊重生命的连续过程与个体生命的不可替代性。

恢复还使我们看到终点与历史的差别。一段运行先中断、后恢复，终点可以与未中断运行一样合格；两段经历却不相同。修复的积极贡献是使某些后续活动重新成为可能，并保存或重建它们需要的通路。已经发生的损伤、错过的时机和应当承担的责任，不会因为终点恢复而自动消失。持续性应当既记录重新取得的能力，也容纳无法撤销的过去。

修复于是包含一种面向过去的忠实和一种面向未来的创造。它要保存足够的连续性，使被恢复的确实是所关切的对象；又可能必须采用不同材料与不同通路，才能重新获得能力。自我维持的成熟形态，不是永远重复相同状态，而是在限制之中维持能够继续变化的组织。\footnote{后续自创生终端保留生命周期、程序来源、编解码、物理动力学与身份桥等输入，详见数学著作第 34 章。正文关于生命能力群的表述是一项哲学组织方案。}

\hypertarget{ux7b2cux516bux7ae0-ux8bb0ux5fc6ux600eux6837ux6210ux4e3aux4e16ux754cux7684ux4e00ux90e8ux5206}{%
\chapter{第八章　记忆怎样成为世界的一部分}\label{ux7b2cux516bux7ae0-ux8bb0ux5fc6ux600eux6837ux6210ux4e3aux4e16ux754cux7684ux4e00ux90e8ux5206}}

设想两次进入外观相同的走廊。此刻能看见的都是同样的两扇门，但入口处先前显示过不同标记：第一次标记说明左门可通行，第二次则说明右门可通行。行动者记得标记，就能够在相同的当前景象下作出不同而有根据的选择。历史已经进入了现在：它通过记忆改变了哪些可能性仍应被认真对待。

记忆因此不只是储存已经结束的过去。它把过去经历转化为对未来的限制，使一个客体能够区分表面相同而行动后果不同的处境。世界中某些重要差别，没有出现在当前观察里，却可以通过观察与行动的连续记录被保留下来。

这种作用说明，信息充分性总要相对于一个问题。若问题只是判断某个动作是否对所有仍可能的情况都符合约束，我们可能只需要知道哪些情况尚未被排除。若问题是比较风险大小、预期收益或代价，就通常还需要概率、评价与时间等更多内容。同一份记忆可以足以支持一种判断，却不足以支持另一种判断。\footnote{核心第 4---5、27 章；数学著作第 37 章的真实历史支持与有限可行性。支持集合足以服务所证明的零未认证占据问题，不据此宣称它足以保持任意概率预测或价值排序。}

这种区别可以在实际过程里检验。先固定同一组允许动作和观察方式，把每段有限行动能够产生的完整观察分布作为预测响应，便能从真实后验推出响应怎样更新。它足以预测这些未来实验，还可在实际过去的概率律上比较其他充分记忆所保存的信息。但已有控制程序可能读取很早以前的记录：压缩预测状态时，不能悄悄删去程序仍在使用的记忆。两段当前景象相同的真实经历，各自都可能有正概率，却要求该程序作出不同动作。保留同一程序的完整历史，才有资格比较它与压缩过程的完整行动记录。\footnote{数学著作第 39.7、39.11---39.12 节。一般有限动作词的真实响应递归与源上的模零最小性，各按其条件陈述；完整历史程序的路径比较保留相同程序和实际历史。没有把不同控制记忆之间的响应相同当作任意行动相同。}

``仍可能''也应当来自所采用的世界模型和实际记录，而不是来自训练资料里见过什么。一个罕见转移没有出现在有限记录里，未必意味着它不可能发生。若把样本中的缺席当作世界中的不可能，就会得到过度乐观的行动空间。真正有意义的记忆压缩，应当说明保留了什么区别，以及排除其他区别的根据是什么。

这使记忆与身份发生联系。一个组织改变了内部编码，只要仍能保存相关历史区别，就可能继续承担原来的预测与行动角色。但这种角色保持并不独自决定它是否还是同一个历史客体。记忆中的信息连续性与承载记忆的历史连续性，需要在具体问题中相接。两个完全相同的记忆副本，可以拥有相同的预测能力，却从复制发生起走向不同的经历。

记忆还有一种容易忽略的社会形态。承诺、档案、语言和共同习惯，使个人无法独自记住的历史继续发挥作用。一个人在陌生城市能够按照规则出行，是因为过去的协调成果被写入公共环境。这种外部记忆扩展了行动能力，也带来了依赖：谁保存记录，谁能修改记录，哪些经历被排除，都可能改变共同世界中的可能性。

所以，信息不是飘浮在物质之外的另一种神秘实体。它是差别在某种组织和任务中被保存、传递并用于行动的方式。它必须有实现，但不能只靠统计材料数量来理解。记忆把一个正在变化的客体与世界的历史相接，使现在不再只是一个孤立瞬间。

\hypertarget{ux7b2cux4e5dux7ae0-ux53efux80fdux6027ux5148ux4e8eux6700ux4f18ux6027}{%
\chapter{第九章　可能性先于最优性}\label{ux7b2cux4e5dux7ae0-ux53efux80fdux6027ux5148ux4e8eux6700ux4f18ux6027}}

当我们问``应该做什么''时，常常立即想到最好、最快或收益最高的行动。但还有一个更早的问题：有哪些行动计划能够真正被实行，并在后续变化中保留所要求的条件？如果没有这样的计划，最优性的计算可能只是在比较不存在的选择。

设想一项有两阶段的任务。第一阶段有两种操作，它们都符合眼前的规范；其中一种把系统带到第二阶段无论怎样操作都无法继续合规的状态，另一种则保留了后续选择。第一步的合格，并不自动等于完整任务的可行。反过来，存在一条完整可行的策略，也不要求每个眼前合格的选择都能够被挽救。

这里的策略不是一条幸运轨迹。它是一种回应方式：在不同的可观察后果发生后，仍能够作出相应选择。若一个动作可能产生两个无法控制的后果，行动者需要为两种正可能后果都保留继续完成任务的方式。只展示其中一个顺利后果，不能说明该动作已经具备完整的可行性。

观察的限制又使问题更严格。行动者可能无法知道隐藏状态究竟是哪一种，因此不能针对不可区分的状态分别选择不同动作。能够实行的，是根据现有观察历史作出同一个选择，并让这一选择对历史尚未排除的状态都满足要求。随后取得新观察，再更新仍然可能的状态集合，继续安排后续行动。

这一思路把可行性理解为一种向未来展开的组织能力。对于已明确给出有限状态、可能转移和时限的模型，可以从终点向前检查：哪些状态集合存在共同合格的动作，且该动作的所有可能观察后继都仍可完成剩余任务。这样得到的``有路可走''，比一句笼统的乐观判断更具体。它依赖准确的可能性描述，也依赖行动者确实拥有相应的观察与行动接口。\footnote{核心第 8、18 章；数学著作第 37 章。该章使用固定有限源、给定观察、时限与明确约束，证明初始支持可行与完整历史策略存在的等价，并给出可达死路的严格例子。}

有人可能据此要求完全排除风险。但必须先说明这里约束的是什么。一个动作可以处于已经完成误差认证的范围内，而结果仍有正的失败概率；它可以维持某种操作合法性，却不保证身体永不受损。可行性只有与明确目标相连，才具有具体内容。把不同层次的``安全''混在一起，反而会让行动者误判保证。

这种有限可行性观念对人的实践也有意义。保存未来选择，并不等于永远推迟决定。有些机会只在承担不可逆承诺之后出现。关键是辨别什么必须提前保留，什么可以合理冒险，以及谁承担失败后果。对自由的理解也应从这里开始：自由既需要选择的意义，也需要实际可实行的选择。

\hypertarget{ux7b2cux5341ux7ae0-ux5b66ux4e60ux5fc5ux987bux56deux5230ux884cux52a8ux6240ux5904ux7684ux4e16ux754c}{%
\chapter{第十章　学习必须回到行动所处的世界}\label{ux7b2cux5341ux7ae0-ux5b66ux4e60ux5fc5ux987bux56deux5230ux884cux52a8ux6240ux5904ux7684ux4e16ux754c}}

一个系统在过去资料上表现良好，我们便愿意让它处理新的任务。然而，资料来自怎样的过程，新行动会访问哪些处境，以及行动怎样由资料决定，是三个不同的问题。只有把它们重新放在同一个生成过程里，``学会了''才获得清楚含义。

设想训练时只尝试过一类温和操作，部署时却因模型的选择大量使用另一类操作。即使前一类操作已经被准确了解，也没有理由把它的平均误差直接当作后一类操作的误差。关于响应的知识，必须到达实际会被使用的地方。覆盖不足不是数据量少的另一种说法；即使重复收集许多熟悉处境，也可能仍未接触决定部署风险的处境。

选择本身也改变了需要考虑的概率。训练记录可能使系统选用不同控制方式，而不同控制方式产生不同未来。评价整个学习过程，应当同时考虑训练怎样发生、训练怎样引起选择，以及选定方式怎样生成部署结果。只分别说明训练和部署的平均性质，可能遗漏它们通过选择形成的联系。\footnote{核心第 16---17、20、27 章；数学著作第 36、38 章。训练实验测度、训练所采行和实际部署占据分别承担不同数学角色。}

这些区分使一种合理的保证成为可能。若同一项训练实验能够以高概率同时校准所有登记的响应，之后根据这些资料选择策略，就可以在同一个校准事件上检验拟合、覆盖和风险。我们无须假装所有可能训练记录都同样好；罕见而误导的记录仍可能存在。保证应当明确承担这些记录的概率，同时利用正常记录上的较强结论。

登记也不要求在每个问题上投入完全相同的观察次数。我们可以在实验前依据关注程度与成本分配不同数量的样本，再共同承担各项响应仍未被准确了解的可能。关键是这些观察确实来自所登记的过程，并且后来的选择仍在这份共同证据的适用范围内。只有在可靠记录上需要用到的结论，也应当在同一批可靠记录上相接；把来自不同筛选条件的好处拼起来，可能没有任何一次实际实验能够同时兑现它们。

这样理解，学习中的统一并不要求训练、表示和行动变成同一件事。我们先辨明源响应的可靠范围，再说明一个新表示保留了哪些差别，最后让所选行动接受这些知识的约束。一份已经共同取得的证据可以支持多个后续选择，无须每换一种描述就把同一份不确定性重新计算一次；但新的描述若遗漏了决定成败的联系，就必须承担这个遗漏。能够复用的是已经取得的理由，不是未经检验的正确性。

这里的校准事件是分析中的真实误差条件，实际操作者未必知道本次训练是否落在其中。可直接检查的经验量，与关于这些经验量接近真实响应的概率保证，应当分别说明。完整的分析恰恰使我们能够在无法直接观察真实误差时，仍对整个实验承担清楚的风险。

本书所依据的数学研究给出了这样一个完整的有限例子：训练是真实指定的随机实验，数据确实改变选择，两种被选控制都有正概率，并产生不同的实际部署分布。整个过程的风险可以从同一联合实验中计算。这个例子的意义在于证明各部分能够真实相接，而非把``选择后仍然成立''作为暗含假设。它同样表明，已经覆盖所有实际访问行，仍不等于失败概率为零。\footnote{数学著作第 38 章的注册选择实例，数值与适用范围见本书附录 C。它证明指定生成实验中的真实联合风险，不从训练边缘律推断任意外部部署机制。}

在认识论上，这推动我们放弃一种过分简单的学习观：学习不是把世界的真相从资料中原样取出，再无损地应用。学习是建立有作用范围的表示，主动选择怎样使用它，并让新的行动后果反过来检验表示。模型的世界与生活的世界需要不断校准。

``不确定''因此是认识进展中的正当位置。它可以指出证据尚不能区分两个解释，也可以指出模型在某处没有得到充分覆盖。若要求每次判断都输出肯定或否定，就可能迫使知识系统把自己的盲区伪装成确定结论。成熟的学习能力，应当包括辨认何时可以行动、何时需要补充信息，以及何时应改变提问方式。

学习也具有伦理后果。由谁决定探索，谁承担试错代价，是否允许受影响者提出模型遗漏，都会影响所谓进步的实际意义。一个能够不断提高预测分数的系统，若从不看见被排除者承担的后果，就可能在技术意义上学习，在共同生活中却越来越盲目。

\part*{\phantomsection\addcontentsline{toc}{part}{第三编　目的、评价与自我修订}第三编　目的、评价与自我修订}

\hypertarget{ux7b2cux5341ux4e00ux7ae0-ux76eeux7684ux4eceux54eaux91ccux8fdbux5165ux81eaux7136}{%
\chapter{第十一章　目的从哪里进入自然}\label{ux7b2cux5341ux4e00ux7ae0-ux76eeux7684ux4eceux54eaux91ccux8fdbux5165ux81eaux7136}}

看见一株植物向某个方向生长，我们可能说它在寻找光；看见一个人绕路，我们可能说他在履行承诺。两种说法都用后来可能实现的结果解释现在的活动，但它们所依赖的根据不同。植物的功能关系需要由生长与维持机制说明，人的承诺则还涉及他对理由的理解。目的语言能够贯通不同领域，前提是我们解释这种贯通怎样发生。

本书主张，目的不必作为自然之外的力量进入世界。它可以存在于有组织的活动中：某些差别被当作成功或失败，系统对偏离作出回应，并且这些回应参与后续行动的形成。在这样的组织里，目的解释指出的是一种真实的方向关系。它并不需要未来事件倒过来推动过去，而是需要现在的机制以某种方式保存、执行或修订对未来的要求。

这个观点的第一层是功能。若某种调节参与组织的维持，我们可以解释它``为了什么''而发挥作用。第二层是形成史：某些组织方式之所以保留下来，与它们过去产生的后果有关。第三层是意向：行动者能够表征结果，比较理由，并据此安排活动。第四层是共同承诺：参与者通过规则把某个结果设为共同任务。这些层次可以相接，但没有理由把它们完全等同。

同一套动力学可以承担不同目的。一个能够前进或后退的装置，可以用来接近某个目标，也可以用来远离它。仅仅知道装置如何运动，无法决定应该采用哪一种评价。即使观察到它长期朝一个方向运动，也仍需区分控制规则、外部约束、偶然初态与行动理由。把所有实际发生的结果都称作系统真正想要的结果，会使目的解释失去区分力。\footnote{原著第 6---8 页；核心第 18、26---27 章；数学著作第 34---35 章。同过程可以具有不同目标；目的机制的识别需要相应忠实性与独立来源。}

目的解释怎样获得比任意归因更强的地位？它需要产生可比较的后果。若我们改变所假定的奖励、信号或维持条件，行为是否会以相应方式改变？若两个目的在过去行为上等价，能否设计使它们产生不同预测的新处境？对于人的行动，还要认真对待当事人的理由陈述，并把陈述与实践之间的关系作为解释内容。主观报告不是不容讨论的最后答案，也不是可以预先忽略的噪声。

反对者可能认为：只要目的依赖评价，目的就只能是主观投射。这个结论过快。评价也可以由真实的组织、制度与历史关系实现。温度偏离设定值之所以触发调节，是机制的一部分；一份契约之所以使行为承担责任，是共同实践的一部分。它们具有现实效力。需要避免的，是把评价有现实来源，误写成任何现实存在都自动决定唯一评价。

自然中的目的因而具有层次性与可修订性。生命形成了基础的利害关系，学习扩大了可表征的未来，语言使理由能够被共同讨论，制度使承诺超越单次相遇。人的目的不是凭空降临，也不是从生存欲望中自动演绎。它在自然和历史中形成，并且能够反过来评价自己的形成条件。

机制与排序之间也可以取得具体联系。在上述寄存器模型中，出生是否实际发生和写入付费都由同一步执行给出。登记价格、折扣和时间后，同一程序在明确费用条件下不劣于原比较类中任一因果策略。实际写入两种值所得的出生均值，加上实际付费均值，在修复机会非零时还能识别目标、费用和故障率，从而确定该登记评价的策略排序。评价合同依然需要理由；机制贡献的是可检验的作用和排序推导。精确均值的数学访问也没有自动解决有限数据学习或伦理正当性。\footnote{数学著作第 39.3---39.5 节；\nolinkurl{RegisterFaultControlModel}、RegisterFaultJointRecord、\nolinkurl{RegisterFaultPurposeInformation}。费用、价格与折扣分别登记；总故障时参数读出退化，并不证明全无差别的目的类不可识别。}

同一机制还可以限制有限资料能够支持怎样的判断。把完整探针批次中的实际出生与付费保留，按预先固定的费用区间进行判断，可以认证某个既有修复程序在登记条件下的未来表现，也会在没有有效信息或费用靠近阈值时留下未知。已经支付的训练费用不能从叙述中删除。更要说明评价的时间：一个修复动作在当前一步付出较多，仍可能使同一程序在较长未来获得正收益。短时负判断不能自动变成永远停止修复的理由。数学证明在这里贡献的是有限而可检查的裁决，评价为何值得采用仍须另有根据。\footnote{数学著作第 39.16 节；\nolinkurl{RegisterFaultFinitePurposeProbe} 与 \nolinkurl{RegisterFaultFinitePurposeInfinite}。有限完整数据的证书、真实无限记录、固定读出下的不可判别和一步／无限时域反转各按其条件陈述。未来最优评价不包括训练总净收益最优，也不等于执行零风险。}

实际付款也限制目的实现。登记的正故障、正费用封闭源中，同一程序从坏态与任意有限储备出发，都有正概率超支；几乎必然破产并未由此证明。\footnote{数学著作第 39.18 节，\nolinkurl{RegisterFaultFeedbackResource}。} 另登记的完整训练协议在原程序、评价和外部供给条件下保留全部故障、出生与费用，使训练前缀非负，并将正修复裁决返回的同一策略接入实际末亲本的随后账本。最优评价从该部署起点计；训练加未来总净收益最优、全部能耗与规范理由各需依据。\footnote{数学著作第 39.24 节，RegisterFaultFundedProbe：完整 \(4N+1\) 次新协议，实际末亲本不同于旧源。}

\hypertarget{ux7b2cux5341ux4e8cux7ae0-ux3c0-ux4e0e-ux3c6ux884cux52a8ux7684ux53ccux91cdux5173ux5207}{%
\chapter{第十二章　Π 与 Φ：行动的双重关切}\label{ux7b2cux5341ux4e8cux7ae0-ux3c0-ux4e0e-ux3c6ux884cux52a8ux7684ux53ccux91cdux5173ux5207}}

一个准备远行的人面临两种不同的困难：他需要足够资源，也需要知道往哪里走。资源再多，错误地图仍可能使他陷入困境；地图再准确，无法获得水和休息也可能使旅程失败。行动通常同时涉及做成事情的条件，以及辨认事情后果的能力。原文的 Π 与 Φ 思想，抓住了这种不能轻易互相替代的双重关切。\footnote{原著第 6 页将 Φ 描述为无知或未适应程度，第 12---15、24---28 页又讨论探索、信息、文化及价值。本章保持这些主题之间的联系，同时明确各指标需要独立定义。}

本书把 Π 用作对资源、实现能力及行动收益的一组问题的提示，把 Φ 用作对信息残差、结构辨识与预测适切性的一组问题的提示。它们的具体含义必须随着模型说明：资源是能量、时间、材料还是某种可交换收益？信息指标测量的是不确定性、预测误差、响应区别还是模型复杂度？不同答案可以建立有用联系，却不能仅因采用相同字母就获得相同量纲。

这种区分能够解释为什么高效活动可能损害长期能力。若把一切可用资源投入眼前产出，系统也许失去监测、学习和维护的余地。反过来，持续收集信息却从不行动，也会消耗支持认识的条件。资源与知识在过程里互相支撑：调查要花时间，准确判断可能节省资源，新的行动又带来新的信息。

但两者不是永远对抗的两股力。一次更好的组织安排，可以同时减少消耗并改善信息；某种损伤也可以同时破坏两者。所谓权衡，应当来自具体可行条件，而不是先验地要求每次改变都沿一条相反方向的直线移动。只有知道哪些改变真正可以实施，我们才能讨论怎样比较代价与收益。

探索还暴露出``信息越多越好''的含混。一个人开始调查陌生领域，可能发现比原来更多的未知。他对自己无知范围的认识增加了，某些不确定性指标却也升高了。这不一定表示学习退步。它可能表示原来的确定感依赖遗漏。信息量、确信程度和理解的准确性，需要分别评价。

同样，真理、道德与文化不能直接被当成一种信息量。一个人可以掌握大量事实而拒绝承担责任，一项制度可以极善于收集资料却压制当事人的发言。若 Φ 被同时用来表示信息、秩序、道德与文明，概念的宽度就会遮蔽它需要解释的联系。更好的做法，是先说明何种知识怎样参与某种有价值的活动，再讨论相应评价。

信息的决策价值还取决于相关性。若一则消息只揭示一枚与当前任务毫无关系的硬币，它可以增加已知事实，却不改善选择。若消息免费、取得它不改变机会而且允许忽略，额外信息不会迫使原有决策变坏；现实的信息活动却往往需要时间、注意和其他资源。关心信息，最终要回到它与实际行动和生活的关系。

双重关切的哲学力量，不在于提前规定自然恰好只有两个驱动力，而在于提醒我们：行动的物质条件与辨认世界的能力彼此相关，却各有无法被另一方取代的作用。一个更强的共同双驱动理论，需要证明不同环境确实共享相应机制，并经受独立的测量和反例。这个开放问题并不削弱双重关切作为分析起点的价值。\footnote{核心第 23、27 章；数学著作第 33---34 章。共同二维秩是特定模型的必要条件；低秩不独自识别物理轴，信息增益也不等于决策价值。}

\hypertarget{ux7b2cux5341ux4e09ux7ae0-ux8bc4ux4ef7ux7684ux5c3aux5ea6ux4e0eux4e0dux53d8ux91cf}{%
\chapter{第十三章　评价的尺度与不变量}\label{ux7b2cux5341ux4e09ux7ae0-ux8bc4ux4ef7ux7684ux5c3aux5ea6ux4e0eux4e0dux53d8ux91cf}}

两位教师可以使用不同分制，却给作业相同的排序。由此可以看出，评价的数值表示与评价本身之间存在差别。知道一个人选了什么，通常不足以知道他在心中使用怎样的精确分数；更不足以知道一个总分由多少独立价值分量组成。

如果总评价由收益和某种代价组合而成，就可能同时调整两项的数值，使它们在总评价中相互抵消。所有最终选择保持不变，收益和代价的解释却已经改变。要识别这两个分量，必须有额外根据，例如独立测量、固定单位、机制约束或能够分别作用于两者的干预。把总目标拟合得很好，并不独自完成分量的识别。\footnote{核心第 18.4、23.2、23.5 节；数学著作第 33、35 章。完全相同数值目标的规范刻画，与仅保持策略排序的目的等价关系需要区别。}

这种表示自由不意味着评价毫无客观内容。它提示我们寻找在合法改写下保持不变的东西：哪些未来仍被优先选择，哪些约束仍不能违反，什么样的干预会改变选择，哪些数值差异由独立计量固定。认识可以先达到一个等价类，再通过新的锚定缩小不确定性。要求每一步都立刻识别唯一数值，反而容易夸大证据。

关系与距离之间也存在类似差别。知道两个状态在某种意义上相同，不等于知道它们偏离时的代价应当如何增长。我们可以用不同刻度记录同样的零误差关系，但相同数字代表的容忍度会改变。若一个组织比较不同层级的损失，必须说明这些刻度怎样校准，否则一处的``小变化''可能是另一处的重大损伤。

这对社会评价具有直接的哲学启发。把健康、时间、尊严、收益和生态影响写入一张总分表，已经作出了一组价值决定。决定有时是必要的，但它们应当能够被说明和争论。不能因为表格形式精确，就把权重当成世界自己给出的事实。一个共同评价制度的合理性，来自它怎样承认相关利益、怎样处理理由以及怎样允许修正。

还有一些价值冲突，并不适合预先压成完整排序。一个人可能认为两种生活各有不可替代的好，而现有理由不足以判定其中一种普遍优越。这种不可比较未必是缺乏理性。它可能反映目标的多元、信息的不充分或作出选择本身的创造性。行动最终需要选择，不代表选择之前就已经存在一个唯一最大值等待计算。

更严格的评价观因此同时承认可比性与它的条件。该计量的要认真计量，该说明的权重要公开说明，该保留的分歧也要保留。数学表示使论证更清楚，而评价的正当性仍需要理由。精确与明智相互支持，各自承担自己的工作。

\hypertarget{ux7b2cux5341ux56dbux7ae0-ux6700ux4f18ux7a33ux6001ux4e0eux5c1aux672aux5199ux5b9aux7684ux672aux6765}{%
\chapter{第十四章　最优、稳态与尚未写定的未来}\label{ux7b2cux5341ux56dbux7ae0-ux6700ux4f18ux7a33ux6001ux4e0eux5c1aux672aux5199ux5b9aux7684ux672aux6765}}

设想一个人每天在工作与休息之间交替。若只看长期平均，他的生活可能呈现稳定比例；若看每一个时刻，他从未停在一个不再变化的状态。这个简单例子说明，长期稳定可以有不同形式。平均结构、反复访问的范围、条件存活者的分布和最终停留的状态，不是同一种事实。\footnote{核心第 10、19、21 章；原著第 6---7、13---15 页关于 GOD／GOA 与演化方向。本文保留长期方向问题，区分不变、平均、循环和吸引。}

原论文用``全局最优方向''与``全局最优吸引子''连接行动与目的。这一设想表达了一个重要愿望：局部活动怎样形成较长期的组织方向？但若把方向预先理解成唯一梯度，把未来预先理解成唯一固定终点，我们会遗漏许多实际可能。不同动作可能同样好，长期行为可能循环，环境可能改变，目标之间也可能发生冲突。

回到第二章的极小装置。它有一个表示运行条件是否满足的可见位，另有一个每步重新取值的内部位。内部位可以改变，也可以保持；只要我们关心的是可见条件和动作后果，就不必逐次分辨这个内部差别。装置有两种动作：维持动作使下一步满足运行条件，中断动作使下一步不满足。这个例子很小，却容许我们认真区分对象描述、行动选择与持续存在。

我们首先看到一种积极的对象性。把四个完整状态按可见位分成两类，可以准确保留每种动作的后果。这样的对象描述有明确的作用：它使我们用较少的区别说明自己正在操纵什么。其准确性不依赖观察者是否喜欢这种分类，而依赖同一类内部的差别确实没有改变所问的后果。若后来要保存内部位的具体内容，就必须扩充描述。一个共同世界可以容许不止一种准确对象描述，因为世界中的问题也不止一种。

让装置从满足运行条件的状态开始，控制方式在每次既往历史之后，都以八分之一的概率选择中断动作。第一次运行后，装置满足与不满足条件的比例就已经稳定：前者为八分之七，后者为八分之一。以后各时刻的分布继续相同。这里的稳定甚至不需要漫长等待。

但如果我们的关切是``从开始以来是否曾经中断''，答案不同。经过两步，曾发生中断的概率是六十四分之十五，高于任一单独时刻的八分之一。有些运行记录先中断，随后恢复，所以只看终点会遗漏已经发生的中断。继续延长时窗，曾经中断的概率不断接近一。一个分布可以从第一步起就稳定，同时每次运行仍不断承受新增的中断机会。

因此，关于持续性的判断要说明我们追踪的是什么。如果关心的是某时刻有多少系统仍在运行，边际分布回答了问题；如果关心一项任务能否不中断地完成，整段历史才是适当对象。后者包含了前者无法保存的经历差别。恢复有真实价值，因为它使后续活动再次成为可能；恢复也不能使已发生的中断从历史中消失。

还可以只考察那些在指定时窗内从未中断的记录。在这些记录中，可见条件始终满足，第一次更新以后的内部位呈现一种稳定的公平分布。这种条件描述也准确。它说明：在继续存活的条件下，系统内部会呈现怎样的结构。可是，进入这个描述之前已经筛去了中断记录。幸存者内部的稳定，不能独自告诉我们原来那套控制方式有多大机会产生幸存者。

这不使条件解释失去意义。相反，它使问题更丰富：存活的机会和存活时的组织结构都值得研究。两种解释可以从同一运行机制严格推出，因而能够互相校验。困难只发生在我们遗忘条件，把关于``继续运行时怎样组织''的回答，当成了``原来的每次运行都能继续''的保证。

最优性又增加了一个维度。假定计分规则奖励中断动作，而维持动作不得分。按照这个明确规则，始终中断得到最高分；若我们同时要求从合格状态出发一直保持运行条件，始终维持才是允许的选择。两种判断都可以计算清楚。它们之间的冲突要求我们审视评价方式和实践约束，不能由``演化一定走向最优''替我们作出决定。最优性需要同时给出可实行的选择、评价方式、时间安排及比较范围；改变任何一项，都可能改变结论。

连``始终维持''的范围也需要说清楚。有一种受限控制只规定：已经满足条件时不得破坏它。对于已经失效的状态，这条规定没有进一步要求。在上述奖励下，它会在合格状态继续维持，在失效状态继续取得中断奖励。于是，从合格起点出发的运行始终合格，整个状态空间却保留了两个不同的长期区域。前一种持续并不需要后者被一个全局终点统一。

这一区别帮助我们理解一种较扎实的方向性。行动可以围绕已经承担的条件形成连贯方向，方向的根据来自当前组织、可实行的回应和所承认的评价。它不必意味着所有可能初态都走向相同命运。对一段具体历史负责，往往首先意味着保存它能够继续承担的关系，而非把整个世界变成只有一种结局的系统。

由此看来，稳定、可行和最优各有积极作用。稳定让我们辨认变化中持续出现的秩序；可行性让我们判断一个承诺是否有实际履行路径；最优性帮助我们在已说明的评价和限制内比较行动。将三者放在同一个过程里，我们既能追问它们怎样相互支持，也能看见它们何时发生冲突。这种统一增加了理解的内容，因为每个概念都必须面对另外几个概念揭示的后果。

至于为什么值得维持某种秩序，还需要进一步的理由。一个安排可能持续很久，却使其中某些参与者承受无法辩护的代价；一种计分方式可能被执行得很准确，却没有记录我们最在意的损失。对目标的检验因此不能停在``它是否被优化''。还应询问谁被纳入评价，哪些后果被排除，以及受影响者能否提出能够改变安排的异议。

这也使意义与持续之间的关系更加具体。对于需要长期履行的承诺，保存身体、记忆、协作和修复通道，可以使一个愿望取得实际兑现的条件。这样的能力值得珍惜。但一件只能完成一次的善意行动、一段已经结束的关系、一项因生命有限而未能完成的事业，也可能具有意义。我们不应把继续运转的能力当作一切意义的必要条件，更不应把一段历史的终结等同于它从未值得存在。

统一自然哲学由此得到一种可以继续发展的方向：在共同世界里，把组织能够怎样保持、行动能够怎样实施、目的怎样得到理由、历史怎样承受后果逐步接合。自然过程为实践提供条件并限制它的可能，人的理由与共同生活使其中一些可能取得值得承担的分量。理解这种联系，可以保留追求统一真理的雄心，同时让每一步认识都有它必须回答的问题。\footnote{数学著作第 21.7、35.7 节给出同一四状态源的完整书面推导：同一策略每个历史条件下以八分之一选择中断，区分其无条件不变律、存活条件分布、曾失败风险及安全受限最优。新数值与完整受限闭环的实例仍按数学正文标明 Lean 范围。本章据此讨论持续和目的；关于可辩护代价、异议与意义的规范判断由本书独立论证。}

\hypertarget{ux7b2cux5341ux4e94ux7ae0-ux80fdux52a8ux6027ux4e0eux76eeux7684ux7684ux81eaux6211ux4feeux8ba2}{%
\chapter{第十五章　能动性与目的的自我修订}\label{ux7b2cux5341ux4e94ux7ae0-ux80fdux52a8ux6027ux4e0eux76eeux7684ux7684ux81eaux6211ux4feeux8ba2}}

一项行动可以有效执行给定目标，而行动者还能够追问目标是否适当，这是两种不同的能力。前者使系统在既定任务中灵活，后者使任务本身成为可以反思的对象。人的能动性很大一部分，来自这两层活动之间的往返。

自我修订并不意味着行动者能站到自身全部历史之外。我们用已有的语言审查理由，用已有的能力想象替代生活，也在社会关系中学会什么值得追问。反思有其条件，但它可以改变这些条件的后续作用。一个人从未承诺为自己选择全部出身，仍可能学会拒绝出身中某些被当作当然的要求。

若把自我理解为一个完全统一的最高控制者，就难以解释内心冲突。不同习惯、情感、时间尺度与社会承诺可能同时提出要求。成熟的能动性不一定要消灭这些差别，而可以形成处理差别的方式：延迟反应、重新描述问题、让较长期的承诺进入当前选择，也允许新理由修正旧承诺。统一自我是不断取得的协调成就。

这种协调不能只以``哪个部分控制了整体''来评价。一个狭窄目标可以非常有效地支配全部活动，却使生活丧失其他重要能力。例如，假设一个人把一切都用于赢得某项排名，他可能高度一致，却无法认真对待友谊或身体的需要。能动性的丰富程度，应当包括识别误指定目标和恢复被排除理由的能力。\footnote{核心第 27 章的反身结构、目标修订与能动性接口；原著第 15---18、24---28 页。关于个人反思与教育的规范意义，由正文提出理由。}

进一步说，知道怎样实现目标，不等于知道为什么要实现目标。技术熟练可以扩大行动空间，但也可能使错误目标更有力量。因而，教育不只应训练执行能力，还应使人能够辨明信息、评价与权力的来源，解释自己的选择，并听取受到选择影响者的回应。

一项反对是：如果每个目标都可以被修订，行动岂不永远无法开始？实际的反思并不要求每次行动都重新审判全部价值。多数时候，我们依赖相对稳定的习惯与承诺；在冲突、失败或重要新证据出现时，才把部分背景带到前景中。理性由此具有层次：它既包含熟练地继续，也包含有根据地中断。

能动性最终是客体与世界之间一种复杂的可修正关系。行动者形成预测，承担某些未来，观察自身行动带来的后果，并在必要时改变表示、程序或目标。它在身体与社会条件中运行。正因为它能够改变条件的后续组织，因果性的存在并不使这种能力变得虚假。

\part*{\phantomsection\addcontentsline{toc}{part}{第四编　组合、尺度与演化}第四编　组合、尺度与演化}

\hypertarget{ux7b2cux5341ux516dux7ae0-ux8bb8ux591aux4e2aux4f53ux4f55ux65f6ux5f62ux6210ux4e00ux4e2aux6574ux4f53}{%
\chapter{第十六章　许多个体何时形成一个整体}\label{ux7b2cux5341ux516dux7ae0-ux8bb8ux591aux4e2aux4f53ux4f55ux65f6ux5f62ux6210ux4e00ux4e2aux6574ux4f53}}

把几个人放进同一房间，尚未形成一个合作组织。若他们开始分配任务、交换必要信息、承接彼此的行动，并维持这种协调，就出现了进一步的问题：新的整体是否具有不能从各自孤立活动直接读出的能力？整体的现实性，要从这种共同活动的后果中说明。

最简单的整体联系，是几个部分共有同一环境。下雨使许多人同时撑伞，但他们并未因此组成共同控制天气的组织。更强的联系来自相互作用：一个人的行动改变另一个人的条件。再进一步，参与者可以共同实现一种有持续性的接口，例如相互校验、共同决策或稳定地产生某项输出。由共同变化到共同组织，需要增加机制内容。\footnote{核心第 15 章的共同性、联合信息、父层形成与组合窗口；数学著作第 32、34 章。本章双卡片为说明联合信息而构造的思想例。}

共同信息也可能超出各部分单独提供的信息。设想一个标记由两张卡片共同编码，只看任何一张都不能确定标记，比较两张是否相同却能读出它。这种新增可读信息没有违背自然规律，也没有神秘地从虚无中产生。它来自两个部分的联合关系。若研究者只测量各部分的孤立表现，就会遗漏这一能力。

这说明整体解释可以有真实内容，而不必诉诸一个脱离部分的额外实体。整体的作用通过部分及其关系实现，但仅罗列零件的独立属性可能不足以预测共同过程。一个乐团的演奏由具体的人和乐器完成，作品的组织却不能被``每个人分别会发声''这一事实充分说明。节奏、分工和相互调整是因果解释的一部分。

与此同时，存在联合能力并不自动赋予整体一个固定身份。一个临时协作可以成功完成任务，然后解散；一个长期组织则还需要维持规则、成员资格、资源和历史连续性。成立一次共同输出、形成可重复机制、出现父层客体，是可以相接的不同阶段。每一步都应当问，它增加了什么可检验的组织内容。

对人的共同生活，这一区分尤其重要。一个群体宣称自己是整体，不能仅凭宣称就取消成员的独立利益。即使共同制度真实存在，仍要讨论哪些能力必须集中，哪些决定适合分散，以及如何保存成员退出、表达和修正的机会。整体性与正当性是两项需要相接的论证。

最好的组合往往既产生新的共同能力，又保存足够的个体性。完全没有约束，协作可能无法维持；消灭所有差别，组织又可能失去学习和纠错的来源。可以把组合看成一个需要维护的窗口：连接足以支持共同任务，差异足以让参与者继续成为有能力的行动者。这个窗口在哪里，应由任务、风险与公共理由共同决定。

整体由此是一种真实成就。它既依赖局部实现，又能通过已经形成的规则与接口改变局部的未来。所谓``自上而下''的作用，在可解释情形中就是这些具体约束的运行；说明约束怎样被实现，比为整体另设一种神秘力量更能保护整体解释的价值。

\hypertarget{ux7b2cux5341ux4e03ux7ae0-ux5c3aux5ea6ux6539ux53d8ux4e86ux4ec0ux4e48}{%
\chapter{第十七章　尺度改变了什么}\label{ux7b2cux5341ux4e03ux7ae0-ux5c3aux5ea6ux6539ux53d8ux4e86ux4ec0ux4e48}}

在地图上看一座城市，街道成为线，建筑成为块，居民的经历几乎完全消失。地图因此不是一张失败的城市照片。它为某类任务保留了重要关系，同时舍去了另一些关系。尺度变化可以使我们获得新的理解，也可以使决定行动成败的差别被隐藏。

粗略描述要成为可靠的有效描述，需要回答一个问题：被合在一起的不同情况，在所关切的后续过程中是否仍能被同样处理？若两条看起来相同的道路，一条允许通行而另一条已经封闭，它们就不能为了导航被无条件合并。压缩的合法性依赖响应与任务，而不只依赖外观相似。\footnote{核心第 7、14---15、20 章；数学著作第 32 章的八类层级关系。预测商、控制商、组合、历史、选择与长期结构各有自己的条件。}

即使压缩保存了预测，也可能没有保存全部控制条件。两个微观状态可以给出相同的某项观察分布，却允许不同的动作，或承担不同代价。要用粗描述直接安排控制，必须把这些差别一并处理。否则，宏观计划可能要求微观现实无法执行的动作。

我们还经常把多种``层级''混在一起。城市由街区组成，是部分与整体关系；把详细地图改成简图，是分辨率关系；追踪城市名称与制度的变迁，是历史延续关系；研究人员在城市之间迁移，是谱系或传播问题；比较城市的公共目标，则进入评价关系。它们共同涉及跨层联系，却并非一种操作的不同名称。统一思想要说明这些联系何时相容，不能提前把它们合为一个万能阶梯。

尺度与时间也需要分开。把地图缩小，并不会使城市已经走过它的未来；长期过程有时可以通过慢变量描述，但这需要证明快慢活动之间的实际关系。一个状态在尺度变化中被省略，不代表它在现实时间里已经消亡。反过来，时间演化改变了组织，也不等于我们仅仅换了一幅更粗的图。

真实系统通常允许近似描述。此时关键不是承诺误差永远不存在，而是说明误差在哪些环节进入、怎样传播以及任务能够容忍多少。一个局部误差若反复放大，就不能因为起初很小而被忽略。若从另一种尺度或单位转移预算，还需要校准比较关系。原有结构的相同，只能保证某些零差别一致，不能自行决定所有非零差别的大小。\footnote{数学著作第 33 章的关系缺陷、双边校准及预算传输。正文关于不同刻度的讨论不被用作物理重整化流已经存在的证明。}

这种观点赋予宏观解释一种有条件的自主性。在已证明的任务范围内，我们可以使用较粗的状态，而不必每次重述所有微观细节。自主性来自可靠的相容关系，不是来自与微观世界脱离。它既允许生物学、心理学与社会科学保留自己的有效概念，也要求这些概念在跨层调用时说明接口。

因此，多尺度统一最值得追求的成果，是明确哪些结构被保持，哪些信息被损失，以及哪些新增机制使宏观规律成立。每一次成功桥接，都让共同世界的一部分联系变得更清楚；它不需要把所有层级提前排进一条已经完成的宇宙流程。

\hypertarget{ux7b2cux5341ux516bux7ae0-ux5171ux540cux751fux6d3bux4e2dux7684ux534fux8c03ux4e0eux51b2ux7a81}{%
\chapter{第十八章　共同生活中的协调与冲突}\label{ux7b2cux5341ux516bux7ae0-ux5171ux540cux751fux6d3bux4e2dux7684ux534fux8c03ux4e0eux51b2ux7a81}}

一家组织可以通过统一安排提高效率，也可以因为统一安排而压制局部信息。成员追求自己的目标，有时能改善整体表现，有时会把代价转移给他人。整体与部分的关系因此不能只用服从或放任来理解。它需要研究目标、信息、权力与后果怎样被组织起来。

最容易产生误解的说法，是``部分应该沿整体的最优方向活动''。这里至少隐藏了两个问题：整体的评价由谁给出？部分的改变是否真的能够改善这一评价？一个方向在某项局部指标上更好，并不保证沿所有局部方向行动会共同提高一个全局目标。不同要求可以互相抵消，也可以形成循环。

设想三种制度安排，每一种都得到一组人的支持，并且各有一种局部比较使它胜过另一种。讨论可能从第一种移到第二种，再移到第三种，最后又回到第一种。若这样的循环真实存在，就不能只给每种安排配一个固定总分，并声称每次局部修正都严格提高总分。我们需要改变协调机制、明确比较条件，或承认存在尚未解决的价值冲突。\footnote{核心第 24 章关于多目标、全局势、父层评价与协调阈值；原著第 8、18---28 页。本章的三方案循环是概念例，用于说明局部改进需要另行连接全局评价。}

因此，协调是一种需要建设的机制，而非所有相互作用自带的终点。信息共享可以减少误解，补偿安排可以重新分配代价，承诺机制可以延长考虑的时间，公开讨论可以修订目标。它们各有条件：信息可能被选择性报告，补偿可能无法覆盖不可逆损失，承诺也可能锁定错误制度。协调的优劣，要连同这些条件一起评价。

多层级组织还需要容纳不同时间尺度。某些局部反馈必须迅速处理，较长期的规则则需要稳定预期。如果所有局部动作都等待最高层决定，组织可能失去响应速度；如果长期约束总被短期信号推翻，共同承诺又无法形成。适当安排不是在快与慢之间简单取平均，而是明确哪类决定由谁在什么时间范围内承担。

一个重要的伦理结论由此浮现：宏观持续不自动给成员提供服从理由。一个制度可能通过把损失长期转嫁给弱势者而维持稳定。若只观察整体是否存续，就会把受损者的处境从评价中删除。共同体的正当性需要说明它如何对待组成它的人，以及这些人能否参与决定维持什么。

这并不支持把一切共同义务还原为即时自愿。共同资源、代际责任和相互依赖会使某些义务超出单次同意。但越是不能轻易退出的关系，就越需要公开理由、适当代表和纠正错误的渠道。整体真实存在，使责任问题更紧迫，而不是使责任问题消失。

统一哲学因此应当把协调理解为维持共同能力的可修正安排。它既反对把冲突视为永恒自然命运，也反对把和谐视为无需机制的必然结果。共同生活的任务，是不断寻找能够同时承受差别与依赖的组织形式。

\hypertarget{ux7b2cux5341ux4e5dux7ae0-ux7e41ux6b96ux9057ux4f20ux4e0eux8c31ux7cfbux7684ux4e2aux4f53ux6027}{%
\chapter{第十九章　繁殖、遗传与谱系的个体性}\label{ux7b2cux5341ux4e5dux7ae0-ux7e41ux6b96ux9057ux4f20ux4e0eux8c31ux7cfbux7684ux4e2aux4f53ux6027}}

一个生命通过修复继续存在，另一个生命通过繁殖出现。两种过程都传递组织，但关系不同。修复力图保持这一生命的连续性，繁殖则产生一个有继承关系的新个体。若把二者都称为``不朽的同一客体''，就会遮蔽个体生命的形成与终结。

谱系的持续具有自己的现实性。某些结构、能力或发展方式通过繁殖机制进入后代，又在后代的活动中被改变和再次传递。这里被追踪的是一条跨个体的继承关系。它可能维持很久，却不使任意两个个体成为同一个生命。家庭中的父母与子女可以深刻地延续彼此，也仍各自拥有生活。\footnote{核心第 25---26 章；数学著作第 34 章关于同一繁殖算子的谱系桥。平均增长、实际灭绝、条件存活和个体身份分别讨论。}

理解遗传，需要区分复制了什么、如何实现以及由什么后果参与保留。相同的抽象模式可以在不同载体中实现，实际传递则依赖具体机制。一个后代在某项预测任务上与亲代等价，未必继承全部功能；一种被广泛传递的模式，也未必提高每个承担者的福祉。传播成功与价值正当性不能互相替代。

选择单位也需要由机制说明。某些变化影响个体留下后代的机会，另一些变化可能只在群体组织或更长期的传递关系中体现。所谓``哪一层被选择''，要考察差异怎样进入繁殖过程，不能仅从我们把什么画成一个整体来决定。对象尺度、组织层级与选择层级可以相互关联，而不必完全重合。

种群的平均变化又不同于每条实际历史。假如一个繁殖过程在平均意义上增长，某些实现仍可能灭绝。平均组成也可能与在存活条件下观察到的组成不同。这样的区别并非技术细枝末节：它决定我们把``会延续''说给谁听，以及把何种风险交给谁承担。

原论文以演化谱系连接生命的短暂与长期组织的延续。这一思路可以保留，同时放弃把演化解释成对一个固定最优终点的必然攀升。竞争者、资源与相互依赖都可能改变选择条件。某种组织今天占优，不代表它已揭示了自然永恒偏爱的形式。谱系的历史既包含累积，也包含损失、分岔和偶然留下的路径。

对人的自我理解而言，谱系意识有双重价值。它使我们看见自己依赖许多未曾亲见的前人，也提醒我们不能把这种继承变成对后代的占有。传递语言、知识与制度，最好同时传递检验和修订它们的能力。真正丰富的继承，给后来者以继续生活的条件，也给他们以成为不同者的空间。

真实出生历史还约束模型应该怎样索引对象。若每个出生都严格进入新一代，一个有限永久个体表就不可能让每个个体总有归一化的后代通道：代数最大的个体已无更年轻的表内子代。这一严格障碍保留了问题，并要求把可重复的类型与一次性的历史个体分开。类型上的平均繁殖公式可以有实际出生来源；它仍不能独自给出每个子代的全局身份或无限谱系的存在。\footnote{数学著作第 39.2 节；\nolinkurl{RegisterRepairHistory.no_finite_normalized_strict_birth} 及 \nolinkurl{actual_birth_price_bridge}。标签历史不是跨时刻并发个体的全局唯一身份。}

标签能够重复，发生的时刻却可以成为另一条历史证据。若真实事务计数在成功、故障和失败时都前进，同一亲本在不同时刻产生的两个子代，即使继承了同一标签，也能由亲本身份、实际计数与标签区分。这条身份来自已发生的记录；擦除时间元数据后，原来的动作、失败与费用分布保持。它仍依赖可靠而可访问的串行记录和登记的初始亲本，不能独自解决并发人口中所有个体的身份分配。\footnote{数学著作第 39.8 节与 \nolinkurl{RegisterTimestampedBirthSource}。同一保留亲本的串行事件身份、成功出生时的父身份解释，与任意并发全局人口身份是不同责任。}

另登记的串行多亲本源从实际查找、修复和新生槽位推导亲本身份，失败仍占时钟；保护程序的复制与随机出生标签分列。冻结旧亲本队列的 Price 只使用其实际子代，空代保留，不把它当成整个末表。\footnote{数学著作第 39.20---39.21 节：串行人口身份与实际子代 Price，空代及原 P5 程序匹配条件保留。} 在另行以自身保护程序确定故障目标的环境中，两类实际子对象分别进入自身维持系统；原固定全局目标的限制继续保留。同一执行与另登记噪声读出的完整律保留全部故障和错误读出，推导实际亲本风险。仅在来源的正确事件且实际出生时，返回的同一个子对象取得维持结论；半噪声时不同亲本仍可有相同读出律。地址、日志与读取权限是这些结论的条件，不从外观猜测祖先，也不推出必然最优或规范正当性。\footnote{数学著作第 39.25 节，\nolinkurl{RegisterOwnProgramJointJournal}。自身程序故障目标与 marked-journal 读取另登记；None 不表示失败交易尝试者。}

\hypertarget{ux7b2cux4e8cux5341ux7ae0-ux521bux65b0ux8106ux5f31ux4e0eux5f00ux653eux7684ux6f14ux5316}{%
\chapter{第二十章　创新、脆弱与开放的演化}\label{ux7b2cux4e8cux5341ux7ae0-ux521bux65b0ux8106ux5f31ux4e0eux5f00ux653eux7684ux6f14ux5316}}

一个组织在熟悉任务中越有效，可能越难应对任务本身的改变。熟练使行动节省资源，探索则可能打断熟练、暴露未知并暂时降低表现。原文以开拓与优化两种倾向理解演化，这为研究创新提供了一个有用起点。\footnote{原著第 12---15 页。核心第 26 章已经区别个体性转变、选择机制与无条件演化最优化；本章继续展开可能空间形成的问题。}

创新有时是在既定选择里找到更好方案，有时则是改变了可以选择什么。一项新协作方式会产生原来不存在的共同能力，一种新表征会使过去难以区分的处境成为不同任务。若把所有创新都画成既定曲面上向高处移动，就容易把选择空间的形成当成早已给定。

因此，需要同时说明变化的两个层次：在固定条件下怎样改善表现，以及条件本身怎样被行动改变。前者适合问哪个动作更好，后者还要问新的动作、关系和评价从哪里来。它们可能在一个历史过程中交替，却没有理由要求每个系统都按照相同生命周期、相同阶段顺序发展。

创新也不天然值得赞美。一个新机制可能扩大合作，也可能提高伤害能力；一种更强的适应可能把代价转嫁给外部环境。变化的新颖性与价值，是不同判断。若仅以能否持续传播来评判创新，我们会错把有效扩张当成充分理由。

脆弱性则提示我们，长期成功可能隐藏条件集中。许多表面独立的活动若依赖同一通路，一次局部中断就可能产生广泛后果。反之，冗余与多样性虽然在当前任务中看似重复，却可能保存未来恢复的机会。它们是否值得维持，仍取决于成本、故障结构和承担损失的方式，不能简单地说越多越好。

一个真正开放的演化观点，应当容纳可逆调整与不可逆转变。恢复既有组织、形成新的整体、产生后代和彻底解体，具有不同历史意义。有时原来的客体消失，部分材料和能力进入新的组织；这种转化是真实的延续关系，但未必是原个体的胜利。正视结束，才能认真讨论怎样保存值得传递的东西。

人的文化为这种开放性增加了反思层面。我们可以意识到自己正在改变选择条件，因而有机会讨论创新的方向、试验的规模与失败后的修复。这样的讨论不会使历史变成可完全设计的工程，但可以减少把任何实际胜出者都解释成必然进步的诱惑。

演化中的希望由此有了更具体的含义。希望是相信仍有能够改善处境的可能，并愿意建设发现和实现这些可能的条件。它不要求世界保证最后成功，却要求我们把影响未来的能力当真。

\part*{\phantomsection\addcontentsline{toc}{part}{第五编　心灵、自由与共同生活}第五编　心灵、自由与共同生活}

\hypertarget{ux7b2cux4e8cux5341ux4e00ux7ae0-ux7ecfux9a8cux4e3aux4f55ux6709ux4e00ux4e2aux4e3bux4f53}{%
\chapter{第二十一章　经验为何有一个主体}\label{ux7b2cux4e8cux5341ux4e00ux7ae0-ux7ecfux9a8cux4e3aux4f55ux6709ux4e00ux4e2aux4e3bux4f53}}

看见颜色、感到疼痛、回忆一次相遇，都不只是关于外部事件的描述。它们对某个主体而言正在发生。我们可以从旁观察一个人的活动，也可以听取他的报告，但这两种接近方式与他正在经历的第一人称经验之间，存在需要认真解释的差别。

本书采取自然主义的起点：主体是世界中的有身体、有历史、有关系的存在者。经验与感知、记忆、行动和身体调节的组织相连。一个主体并非先完整地存在于内心，再通过通道收集外部图像；他在与世界打交道的过程中形成能力、视角和自我理解。这个起点使心灵能够被纳入共同自然，而不必被放在因果世界之外。\footnote{原著第 15---18 页；核心第 27.6 节区分功能、意向与规范意义。本章的主体自然主义是哲学立场，未将神经组织模型等同于完整的现象意识解释。}

从组织角度看，经验具有某种统一性。声音、触感、意图和回忆并不是永远互不相干的碎片，它们可以共同影响一次行动。与此同时，这种统一并不要求全部内容在任何时刻都被一个中央观察者看见。很多活动在注意之外展开，某些冲突直到行动中断才变得显著。主体的统一，可以是多个过程之间持续取得的整合关系。

但``功能被整合''与``为什么有体验''之间，还需要进一步论证。一个模型可以准确描述信息怎样传递、哪些内容可以报告、行为怎样改变，却仍未单独回答痛苦为何作为痛苦被经历。承认这个问题，能够保护自然主义的解释责任：不能把尚待解释的主观性，仅以``复杂网络''重新命名后就认为已经完成。

一种可能的反对是：既然组织说明尚未解决全部经验问题，是否应当把意识放回自然之外？这同样走得过快。解释尚不完整，不会自动提供另一种实体存在的理由。我们可以逐步寻找身体、过程、报告与经验之间更强的关联，检验不同组织假说，并保留第一人称材料作为研究约束。自然主义的力量在于它允许这些证据进入一个共同问题，而不在于提前否认问题。

相反，把所有有反馈的客体都称为有意识，也不足以解决困难。一个装置可以维持变量，一条谱系可以保留信息，却未必因此拥有我们谈论体验时所指的现象。客体性、功能性主体、反思自我和现象经验，需要通过具体论证关联。概念有连续的形成条件，不意味着所有区别都已消失。

主体性还含有不可替代的处境。我的身体在这里受伤，我的记忆把这次遭遇与某段历史相接，我的行动从这一位置改变世界。别人可以理解、帮助并解释我，却不能通过把我的经历写入一个总模型而使我成为多余。第一人称地位在共同世界中有位置，正因为共同世界容纳着不同的真实承受者。

这样的理解为心灵与物质之间提供了一条开放道路。物质组织使经验生活成为可能；经验中的理由又通过行动参与世界。两种说明可以相互深化。自然的统一，不必以取消主体为代价；主体的真实性，也不需要通过否认身体和关系来维护。

\hypertarget{ux7b2cux4e8cux5341ux4e8cux7ae0-ux6ce8ux610fux60c5ux611fux4e0eux6709ux5206ux91cfux7684ux4e16ux754c}{%
\chapter{第二十二章　注意、情感与有分量的世界}\label{ux7b2cux4e8cux5341ux4e8cux7ae0-ux6ce8ux610fux60c5ux611fux4e0eux6709ux5206ux91cfux7684ux4e16ux754c}}

如果所有进入感官的变化都同样重要，行动就难以组织。注意使某些内容成为前景，使另一些内容留在背景；情感则使世界呈现紧迫、可亲、危险、失落或值得期待的不同分量。它们参与构成一个行动者实际面对的世界，而不只是给已经完成的推理附加色彩。

注意的选择需要历史。熟练者在同一场景中看见的关键线索，可能是初学者尚未辨明的差别。关心某个人，也会使我们比过去更容易注意到他的需要。目标、记忆与感受共同塑造显著性，而显著性的变化又会影响新的目标。这样的回返关系，使注意成为能动性的一部分。

情感同样可以含有认识作用。担忧使尚未发生的损失进入考虑，羞愧让自己在他人眼中的行为成为问题，愤怒可能使被忽略的不公显现。但情感也可能误认对象、夸大威胁或固守旧经验。因此，认真对待情感既不要求把它当作绝对权威，也不要求先把它排除，才开始理性思考。理性可以通过情感认识利害，又通过新的信息和理由修订情感所指向的解释。

这里需要区别调节与压制。一个人能够延迟即时反应，给其他理由进入选择的机会，这可以扩大自由；若他只是学会对所有不适保持沉默，也可能失去辨认损伤的能力。调节的价值要看它保存了什么，而不能只看外部行为是否平稳。稳定的表情可以伴随丰富的自主，也可以伴随不得表达的处境。

自我理解有助于这种调节。行动者能够发现：某个反复出现的反应来自过去经验，它曾经有效，却未必适合当前环境。这种理解并不使反应立即消失，但可以为改变创造条件。新的习惯、环境安排与他人的支持，也能参与形成不同反应。意志的工作因此常常是组织条件，而不是由一个无条件强大的内部命令压倒身体。

同理心则把这种理解延伸到他人。我们通过他人的表达、处境及自身类似经历，尝试把世界放到另一个视角下理解。然而，想象别人的感觉并不等于已经知道。对方的纠正具有重要地位。真正的理解需要保留听取异议的通道，尤其当我们的模型已经非常自信时。

原论文把注意、情感、记忆、意志和同理心置于动态组织中，这个方向值得发展。具体的神经机制需要相应证据，概念层面的共同结构则可以先帮助我们提出问题：哪些内容被整合，哪些反馈受到控制，哪些信息被忽略，哪些目标可以修订。\footnote{原著第 15---18 页的注意、情感、记忆和同理心主题。本文不采用未经独立确认的神经数值、临床诊断指标或具体疗效主张。}

有分量的世界不是客观世界的虚假版本。它是一个有限主体从自己的处境与关切进入现实的方式。问题在于，这种进入方式能否继续接受世界与他人的修正。认识的成熟，包括让自己的显著性结构变得更敏感，而不只是让既有关注变得更强。

\hypertarget{ux7b2cux4e8cux5341ux4e09ux7ae0-ux81eaux7531ux4e0dux662fux6ca1ux6709ux6761ux4ef6}{%
\chapter{第二十三章　自由不是没有条件}\label{ux7b2cux4e8cux5341ux4e09ux7ae0-ux81eaux7531ux4e0dux662fux6ca1ux6709ux6761ux4ef6}}

一个人在威胁下说``我同意''，与他经过理解和讨论后说同样的话，具有不同的行动意义。两次发声都发生在自然因果过程中，但自由程度显然不同。自由的解释因此不能只问行为有没有原因，还要问原因以怎样的方式经过行动者。

本书把自由理解为一种有条件而真实的能力：行动者能够理解相关处境，形成和修订意图，依据理由在实际可行的选择之间行动，并拥有适当的机会抵抗强制。这种自由可以逐步发展，也可能被削弱。它不需要一个脱离身体与历史的起点，却需要行动者在过程中的作用不能被抹去。

一种反对认为：只要行动受到性格、环境和过去经验影响，就不是自由。若接受这个要求，理解与学习也会成为自由的敌人，因为理由越能影响一个人，他的行动就越有原因。更合理的区分，是哪些影响发展了行动者判断和选择的能力，哪些影响绕过或破坏了这种能力。被说服与被操纵并非因为一个有原因、一个无原因而不同，而是因为它们给予理解、纠正和拒绝的地位不同。

另一种反对担心，自由若有程度，便可以按某个统一能力分数给人排等级。这个推论也不成立。评价特定行动的控制能力，目的在于合理安排支持与责任，不是决定谁拥有作为人的基本地位。一个人在某项任务上能力受限，仍然是其生活后果的真实承受者。尊重并不以达到最高控制复杂度为条件。

责任判断需要考虑知识、控制、选择机会和可合理预见的后果。一个人若无法取得关键信息，与明知后果却隐瞒信息，不应被同样评价。制度若把不可完成的任务强加给参与者，也不能在失败后仅以``他作出了选择''卸除自己的责任。自由的实际空间由个人能力和公共条件共同构成。\footnote{原著第 24---28 页；核心第 18、27 章的可实行策略与能动性分量。自由与责任的公共理由属于本书的独立规范论证。}

这也使反思与行动可能性发生联系。一个人可以想象许多选项，但若每个选项都被同一种强制堵住，想象本身不足以构成充分自由。反过来，有许多形式选项，若无法理解其差别，选择也可能只是随机移动。自由需要认识的能力与可行的机会相接。

个人自由和共同规则之间并非只有对立。承诺可信、公共安全、表达权利和可预测的程序，可以扩大每个人能够实现的计划。规则也可能被用于支配。判断关键是规则如何产生、作用怎样分配，以及受规则影响的人能否对它提出有后果的异议。

自然主义由此能够认真对待自由。一个能够通过理由、记忆与共同制度改变未来的客体，拥有真实的因果能力。让这种能力得到发展，正是一项现实而非超自然的任务。自由不是世界停止对人发生作用，而是人在与世界的联系中逐步取得可以理解、承担并修正自身行动的位置。

\hypertarget{ux7b2cux4e8cux5341ux56dbux7ae0-ux610fux4e49ux6765ux81eaux600eux6837ux7684ux627fux8bfa}{%
\chapter{第二十四章　意义来自怎样的承诺}\label{ux7b2cux4e8cux5341ux56dbux7ae0-ux610fux4e49ux6765ux81eaux600eux6837ux7684ux627fux8bfa}}

当一个人说生活失去意义，他通常不只是报告效率下降。他可能仍能完成任务，却不知道这些任务为什么值得完成；也可能真诚关心某件事，却被现实条件阻断了参与它的机会。意义问题涉及主体、行动与价值之间的联系。

本书把意义理解为这样一种关系：行动者能够把自己的活动理解为参与某种值得关切的事物，并使这种参与在生活历史中取得位置。这里包含经验的投入、实际的联系和可讨论的价值。三者可以相互支撑，也可能分离。一个人觉得充实时，仍可以追问他所投入的事业是否伤害他人；一项确有价值的工作，也可能因为处境恶劣而无法给承担者带来完整感。

意义不必等待一个唯一的人生终点来赋予。照顾一个具体的人、完成一项诚实研究、保存一段重要记忆、修复一种受损关系，都可以在有限历史中具有分量。它们之所以有意义，不要求整个宇宙事先为此规定一个总目的，而是因为有真实的承受者、现实的关系与可以说明的价值。

这不等于意义完全由个人任意创造。一个人能够宣称任何事情对自己有意义，但关于意义的公共讨论仍可询问：这种活动实际连接了什么？它是否建立在欺骗上？它使谁能够更好生活，又让谁承担不可承认的代价？意义中的主观投入十分重要，但投入的强烈程度不能独自证明对象值得投入。

原论文把个人与世界的动态协调置于意义中心。这一洞见可以进一步展开：个人的计划依赖身体、他人和公共环境，意义经常从参与超越孤立自我的活动而来。然而，协调不能被理解成无条件顺应现存整体。一个拒绝不公命令的人，可能与所在组织发生冲突，却正因承担了有理由的反对而获得深刻意义。社会稳定的方向并不总是值得认同的方向。\footnote{原著第 24---28、40---41 页；核心第 27.6 节。本文区分意义的体验、实际参与及值得追求的根据，并用正当异议检验``与整体对齐''的表达。}

对于持续履行型的承诺，意义的实行需要相应条件。照护需要有人在场，研究需要能够延续的记忆与交流，共同计划需要可以信任的回应。保存这些条件，使愿望逐步成为可承担的实践。修复的价值也在这里：它有时能重新打开履行的机会。但这些是特定承诺的实行条件，并非一切意义都必须具有的条件。一次正当的拒绝、一项无法重复的善意行动，也可以在其发生时取得完整的分量。

意义还需要容纳无法完成的承诺。一个人的事业可能中断，关系可能因死亡而结束，能力也可能改变。若把意义等同于不断接近可测的成功终点，这些生命阶段就会被轻率地判为失败。实际上，忠实地关切、适当地哀悼、把尚未完成的工作交给后来者，也可以具有不可替代的意义。

有限性使这种意义更具体。正因为时间有限，承诺意味着放弃某些可能并认真进入另一些可能；正因为人会受伤，照护不是无差别交换；正因为历史不能完全重来，诚实与修复才有分量。生命的意义不必通过抹去有限性来保存。

因此，意义是需要维持又能够修订的联系。个人既可以通过承诺形成连贯，也可以在发现更好理由时改变承诺。共同体的责任，是为这种形成和修订提供空间：让人能够贡献、被听见、获得支持，并在必要时重新开始。一个统一世界观若能帮助人更准确地理解这些联系，就已经对意义作出了实质贡献。

\hypertarget{ux7b2cux4e8cux5341ux4e94ux7ae0-ux9053ux5fb7ux9700ux8981ux53efux4ee5ux76f8ux4e92ux7ed9ux51faux7684ux7406ux7531}{%
\chapter{第二十五章　道德需要可以相互给出的理由}\label{ux7b2cux4e8cux5341ux4e94ux7ae0-ux9053ux5fb7ux9700ux8981ux53efux4ee5ux76f8ux4e92ux7ed9ux51faux7684ux7406ux7531}}

一个组织可以凭借恐惧维持秩序，也可以通过公平合作维持秩序。若只考察哪种方式使组织持续，我们尚未回答哪种方式值得人们接受。自然过程提供了利害与依赖，伦理学则需要说明，为什么某种利害应当成为他人的行动理由。

本书采取三项相互关联的规范立场：真实的痛苦与能力损失值得被认真对待；相似处境中的人有权要求相近理由，而不能仅因身份不同被任意排除；共同规则应尽可能允许受影响者参与解释、质疑和修订。这些立场得到身体的脆弱性、相互依赖以及公共说理实践的支持，但它们仍需通过规范论证被辩护，不能伪装成从生物存续直接推出的定理。

为什么这些立场值得接受？首先，它们要求我们在第一人称与第三人称之间保持可说明的一致性。若我认为自己的痛苦不能仅因别人获益就被忽略，我就需要解释，为什么相近的痛苦发生在别人身上便失去全部理由地位。其次，它们能够支持一种较可靠的共同生活：每个人都不必完全依赖他人的偶然善意，才能使自己的基本处境进入决定。最后，它们给修正错误留下空间，因为受损者可以提出我们尚未看到的事实与理由。

这套论证并不假定所有利益总能完全兼容。资源有限、承诺冲突、紧急处境，都可能迫使人承担损失。正因为如此，需要区分必要负担与任意牺牲，比较替代方案，说明决定程序，并认真承担补偿和修复。道德推理不只寻找一个总结果，还关心人以怎样的地位参与结果的形成。

一个纯粹利己者可以拒绝相互承认。对此，不能假装只用一条自然规律就能迫使他接受全部义务。我们能够指出，他对自己的要求与对他人的排除需要额外理由；能够说明这种排除怎样破坏公共信任；也能建设防止支配的制度。伦理的力量包括说理、教育与制度实践，而不只是找到一个无法逃避的演绎开端。

关怀与公正因此需要结合。关怀使我们看见具体的人，公正提醒我们不能只对熟悉的人敏感。一个由同情心驱动的群体可能仍忽视远方或沉默者；一个只处理抽象份额的程序也可能错过实际困难。制度应当把普遍理由与具体处境连接起来，而不是要求人只能选择冷漠的平等或偏私的温暖。

自然主义在这里能够发挥积极作用。它帮助我们理解能力怎样受身体与环境影响，制度怎样把某些人置于脆弱位置，以及照护和教育怎样扩展行动能力。伦理学则进一步讨论哪些扩展值得共同承担，哪些代价不能轻易转嫁。事实研究与规范判断通过明确理由相接，两者都能修正对方的问题设定。\footnote{原著第 24---28 页的自然主义伦理问题；核心第 18、24、27 章明确保留评价桥与主体间约束。本书提出的关怀、相互承认和公共理由原则，由正文辩护。}

道德的普遍性也不意味着所有生活必须相同。我们可以要求相互尊重和基本理由的一致，同时允许不同的人追求不同的善。共同世界值得维护，正因为其中存在多种有分量的生活。一个更成熟的统一哲学，应该使这种多样性得到组织，而非以统一之名取消它。

\hypertarget{ux7b2cux4e8cux5341ux516dux7ae0-ux5236ux5ea6ux5e02ux573aux4e0eux96c6ux4f53ux7406ux6027ux7684ux6761ux4ef6}{%
\chapter{第二十六章　制度、市场与集体理性的条件}\label{ux7b2cux4e8cux5341ux516dux7ae0-ux5236ux5ea6ux5e02ux573aux4e0eux96c6ux4f53ux7406ux6027ux7684ux6761ux4ef6}}

一项制度把很多次相遇连接起来。承诺可以跨越时间，规则可以跨越个人在场，某种角色可以在承担者更换后继续发挥作用。制度因而是一种真实的组织：它通过信息、资源、资格与预期改变行动的可能性。把制度看作客体，有助于研究它怎样形成、怎样维持以及怎样修复。\footnote{原著第 8、18---28 页关于层级和社会客体；核心第 15、24、29 章。本文将制度的组织性与制度的正当性作为相连而不同的问题。}

制度的客体性同时依赖参与者的实践。写在纸上的规则需要被理解、执行和争论；统计指标需要测量和使用；角色需要有人承接。这使制度既具有超出单个人的持续性，又始终可以追问其具体实现。所谓``制度要求''，不能永远作为停止解释的最后一句话。

市场可以作为一个重要例子。价格、承诺、产权、信息与交换规则，把许多分散行动连接起来。在合适条件下，局部信号能够引导资源流向新的需求，并使参与者无需掌握全部世界信息就开展合作。但这些条件本身是制度组织的一部分：哪些成本进入价格，谁掌握信息，何种承诺可信，以及谁能够参与交换，都会影响结果。

因此，从个体逐利到共同利益，需要机制性的论证。一个人的收益可以来自为他人提供真实价值，也可以来自把损失转给没有谈判力量的人。若后果没有进入交易评价，局部最优就可能累积成公共损伤。原文对市场泡沫、集中与危机的关注，可以在这里获得更准确的哲学位置：研究局部反馈怎样通过制度连接，而不预先假定它们必定通向集体最优。\footnote{原著第 18---24 页。原稿中的泡沫、财富集中与危机是本章的问题来源；本文不把特定分布律、风险指标或政策效果当作已由 UEOT 确认的普遍事实。}

效率与公正同样需要分开。设想两个人分配总量固定的一种资源，双方都把更多资源视为改进，全部资源已经被其中一人持有。在只改变份额的比较中，改善另一人的处境必须减少持有者的份额。这个安排可以满足``无法在不损害任何人的情况下再改善''的条件，却仍需要回答分配是否公正。效率判断的比较规则，不能独自决定权利和初始分配的正当性。

财富积累和网络集中也应当以具体机制研究。收益能够再投资，信息优势能够带来新的机会，规则能够强化既有地位，这些都可能形成累积过程。但某种反馈存在，并不自动给出所有社会相同的分布规律，更不自动决定一种唯一改革措施。要比较解释，需要看到转移通道、时间范围、替代机制和受影响群体。

系统风险经常来自不同时间尺度的错配。快速的局部反应可以改善单个参与者的处境，也可能共同撤去整体依赖的支撑；缓慢的制度修正可以稳定预期，也可能来不及处理已经发生的变化。把网络画出来，能够指出联系，却还需要知道每条联系承载什么义务、怎样传播损失以及有哪些缓冲。一个漂亮的总体指标不能替代这些问题。

从这些认识出发，公共治理最基本的任务，是使共同后果能够回到共同决定之中。公开信息、明确责任、适当分散权力、为异议提供通道、保留局部失败后的恢复能力，都可以作为制度设计的问题。其具体实现应由经验和公共讨论决定。哲学能够澄清评价目标，不能仅凭一个统一框架替所有社会计算出同一份政策。

集体理性也不等于所有人拥有同一个头脑。它可以通过可靠的分工、彼此校验和可纠正的程序形成。真正值得追求的整体，是一种让分散知识获得作用、让受损者能够发声、让公共错误能够被承认的组织。这样的整体会有分歧，但分歧本身可以成为其认识能力的一部分。

\hypertarget{ux7b2cux4e8cux5341ux4e03ux7ae0-ux6587ux660eux6280ux672fux4e0eux6bd4ux6211ux4eecux66f4ux5927ux7684ux4e16ux754c}{%
\chapter{第二十七章　文明、技术与比我们更大的世界}\label{ux7b2cux4e8cux5341ux4e03ux7ae0-ux6587ux660eux6280ux672fux4e0eux6bd4ux6211ux4eecux66f4ux5927ux7684ux4e16ux754c}}

个人的行动依赖语言、工具、制度与自然条件，这些条件大多不是个人独自创造的。文明可以被理解为一种跨世代维持和改造行动条件的过程：保存知识，形成公共规则，建设基础设施，同时也制造新的依赖和风险。它把许多生命连接进更长的历史。

原论文描绘了一个朝向更高信息组织与更广整体协调发展的文明图景。这种图景能够激发责任感，但文明进步不能只由信息数量、资源效率或系统规模衡量。更强的组织能力既能改善照护，也能扩大监视；更快的传播既能分享知识，也能扩散误导。评价进步，要问能力怎样分配，哪些生活得到改善，以及新的力量是否能够被共同约束。\footnote{原著第 24---28、40---41 页关于主体---世界复合与文明方向。本文把文明成熟作为需要价值辩护的方向，未假定信息／资源比例随历史必然单调提高。}

技术使这一问题更加鲜明。一个自动化系统可以在给定任务中表现出复杂的预测、选择与自我调整，这些能力值得认真研究。但把任务执行得更好，并不自动给出适当的任务，也不自动决定谁有权部署它。技术的目的与责任需要通过实际制度安排解释，不能归给一项仿佛天然中立的效率趋势。

对于具有复杂学习能力的系统，理解、能动性和经验地位也应分别讨论。我们能够研究它保存什么历史、怎样选择行动、怎样响应理由，而关于是否存在主观体验，还需相应论证。谨慎的区分有助于避免两种实践失误：因为尚不确定体验而忽视真实影响，或因为能力惊人而把所有责任交给一个未经说明的``自主主体''。

自然环境则使文明看到自身更深的开放性。一个社会可以在账面上增长，同时耗损它所依赖的恢复条件。把外部世界只当作可替换输入，会遗漏许多变化的时间尺度和不可逆性。生态思考要求追问，哪些依赖不由当前交易直接记录，哪些后果会跨越地域与世代。

这并不要求把地球想象成具有唯一意志的超级个人。我们可以承认生态系统的真实联系，而不为它假定一个可以直接读出的总目的。人对非人生命及自然过程的责任，可以来自多种理由：它们支撑人类生活，它们包含值得关切的生命与经验，它们具有无法由即时利益充分衡量的历史和关系。具体伦理立场可以争论，现实依赖却要求这种争论不能永远被搁置。

代际责任尤其显示了共同体的开放边界。后来者无法参与今天的全部决定，却要生活在这些决定留下的条件里。我们无法替他们决定完整的善生活，但可以尽量保存他们形成、修订并实行计划的能力。维护可恢复性、知识的可检查性以及公共制度的可修正性，因而具有超越当下偏好的理由。

文明的成熟可以据此理解为一种反身能力：不仅增加改变世界的力量，还能够看见力量的前提、分配与后果。一个更统一的世界观，应使我们更准确地感受这些联系，而不是赋予文明一个永远扩张便自动正当的使命。

比我们更大的世界，无须通过一项宇宙命令才值得认真对待。它已经通过依赖、相遇、脆弱与共同历史进入生活。承认这种事实，同时为关切提供公开理由，是自然哲学与伦理学可以共同推进的工作。

\part*{\phantomsection\addcontentsline{toc}{part}{第六编　统一的方法与真理的进步}第六编　统一的方法与真理的进步}

\hypertarget{ux7b2cux4e8cux5341ux516bux7ae0-ux8ba9ux89e3ux91caux627fux62c5ux5931ux8d25ux7684ux53efux80fd}{%
\chapter{第二十八章　让解释承担失败的可能}\label{ux7b2cux4e8cux5341ux516bux7ae0-ux8ba9ux89e3ux91caux627fux62c5ux5931ux8d25ux7684ux53efux80fd}}

一个框架如果能够为任何结果找到事后说法，就很难告诉我们原来不知道什么。解释的力量来自它在结果出现之前承担了某些可能的失败：它指出哪些差别重要，哪些后果应当出现，以及什么观察会促使我们改写说明。

原论文提出从客体划分、结构剖面、双重关切、方向与长期组织到综合诊断的分析路线。这条路线可以改造成一种有纪律的提问方式。先说明我们追踪什么，再辨明实际接口和维持依赖，然后提出评价与机制解释，最后用独立结果检查这些解释之间的联系。这样，统一概念不会替具体工作省去责任，而会使责任更容易看见。\footnote{原著第 9---11、37---40 页；核心第 16---17、28---31 章，以及科学接口资料关于独立认证、反例与不确定结论的区分。}

一个解释首先要保护问题本身的稳定。如果任务已经失败，研究者事后改称那一结果从来不属于客体自己的未来，就会让客体边界成为逃避反驳的工具。边界可以修订，但修订应当说明它改变了原来哪项主张。持续使用一个名称，不能抹去对象、协议或评价已经变化的事实。

第二，需要区别概念关系与定量尺度。说某种压力偏离了组织的方向，还不足以定义一个普遍的偏离分数。方向本身怎样取得，零压力时怎样处理，完全相反的方向是否被指标错误归为同轴，都需要检查。若把不同单位的资源和风险合成比值，还应说明校准方法。原文中的诊断指标可以成为待研究的工具，但工具的名称不能代替其有效性。\footnote{原著第 9---11 页的偏好强度与偏离指标。本章提出的是指标定义与校准的检查问题，没有为其提供新的普遍经验阈值。}

第三，要把模型成功与现实来源相接。一份关于已知过程的证明，可以说明如果世界满足这些条件，某个结论必然成立；一项模拟可以显示这些条件在一个具体构造里共同实现；现实实验则要检验外部过程是否确实具有所需性质。三种工作能够相互支持，分别回答不同问题。由研究者生成的世界，不会因为生成得复杂而自动成为现实世界的确认。

在有限数据情形中，事先固定的比较、保留全部尝试、说明样本如何取得，以及独立复核失败，能够减少选择性解释。这里的重点不是仪式化地增加表格，而是使别人能够知道：研究者看见过什么，哪些结果未被采用，以及同一证据在多少项结论中被使用。一个共同校准事件被多次调用，不应被误算为多次独立成功。

失败也需要准确归因。仪器出了问题，不直接推翻所有理论；模型的应用前提不成立，不等于演绎证明出现矛盾；评价桥没有识别真实目的，也不意味着被研究客体不存在。但这些区分必须对批评开放，不能由理论提出者单方面把所有不利结果送到外围。若多个合理适配都系统性失败，就需要重新审视更核心的假设。

对统一研究尤其有价值的是具有区分力的对照。若几种理论都预测会出现某种循环，仅发现循环就不能单独支持其中一种。需要比较它们对干预、时间尺度或新环境的不同预测。实验的决定性来自竞争解释被怎样区分，而不是来自使用了多宏大的观测对象。

解释承担失败的可能，也意味着承认``不确定''有时是最好的当前结论。确定性必须由证据取得，不能作为一种态度强加给世界。能够准确定位未知的理论，已经帮助后来者减少了寻找方向的成本。

\hypertarget{ux7b2cux4e8cux5341ux4e5dux7ae0-ux5b87ux5b99ux7edfux4e00ux4f5cux4e3aux89e3ux91caux7684ux8fdcux666f}{%
\chapter{第二十九章　宇宙统一作为解释的远景}\label{ux7b2cux4e8cux5341ux4e5dux7ae0-ux5b87ux5b99ux7edfux4e00ux4f5cux4e3aux89e3ux91caux7684ux8fdcux666f}}

人类能够在星空、生命与自身心灵之间感到联系，这种体验并不自动是一项科学定理，却提出了真实的哲学问题：世界的多样性是否具有更深的共同条件？拓扑、信息、动力学与组织能否在不同尺度上形成相容的解释？原始 UEOT 论文把这些问题推进到基础物理与宇宙整体，显示了它最广阔的企图。\footnote{原著第 29---39 页的基础物理与宇宙学猜想。本文仅讨论它们对统一自然哲学提出的解释责任；没有把这些猜想纳入本轮核心数学结果的完成范围。}

本书保留这种远景。寻求统一可以促使我们发现重复出现的结构，追问某些看似独立的原则是否来自共同机制，并尝试从更少的假设理解更多现象。哲学能够为这项工作提供概念方向，也能够要求它说明何种成果才算真正的推导。

一个共同词汇还不是共同理论。若把不确定性、生命适应和社会无知都叫作``信息''，我们只是提出了一个比较邀请。要形成统一说明，需要定义各自的信息对象，说明它们怎样对应，以及一个领域的定理为何能够传到另一个领域。同样，把轨迹循环、因果依赖和空间拓扑都用 Ω 表示，可以保留家族相似性，却不能省去这些数学与物理对象之间的桥。

原论文提出了三组基础物理猜想：从组织和信息理解量子结构，从拓扑组织解释粒子属性，从集体信息关系理解时空与引力。它们在本书中作为统一研究的方向保留，而非已经完成的自然图景。若从统计区分走向时空几何，需要说明度量、因果结构与动力学为何取得所需形式；若从拓扑分类走向粒子理论，需要说明可观察量、相互作用及数值预测的对应；若用尺度变化连接不同理论，需要给出实际相容的过程。每一项都是能够取得独立进展的研究责任。

尤其应当区别``我们可以构造某种描述''与``现实选择了这种描述所对应的机制''。数学通常允许许多对象和多种表示。物理解释还要说明为何相应条件适用，参数怎样确定，替代模型怎样被排除。统一的美感可以引导寻找，选择则需要由推导与观察共同完成。

原文最后提出，客体组织与双重关切可能是一个更深作用原则的不同表现。作为形而上学猜想，这表达了世界的秩序也许具有共同生成条件。作为数学物理研究，它则需要一个明确的作用对象、合法变化范围、动力学方程与现实解释。只有这些内容逐步形成，统一作用量才从名称成为有后果的理论。\footnote{原著第 41 页的终极统一作用原则猜想。本文保留其作为研究远景的地位，不以其存在作为已经完成的数学或物理前提。}

宇宙整体的目的还面临额外问题。我们可以讨论一个生命怎样维持与环境的交换，因为它有可区分的边界和条件；把同样说法移到全部宇宙，需要解释这里的环境、维持与成功究竟指什么。即使宇宙具有整体结构，也不因此自动成为具有意向的个人。整体组织、生命式自维持和体验中的目的，需要分别论证。

这并不使宇宙在意义上贫乏。人能够认识部分自然，生命能够形成关切，社会能够共同追问真实，这些活动本身已经发生在宇宙中。意义的存在不要求每一部分都拥有同一种意义，也不要求宇宙整体替我们作出伦理选择。一个容纳真实主体与共同价值实践的自然世界，已经比``只有无关事件堆积''的图景丰富。

因而，本书对终极统一的态度是有方向的开放。值得寻找使不同解释相容的更深结构，值得提出大胆猜想，也值得在必要时放弃漂亮的桥接。自然的统一如果真实存在，就应当能够承受我们逐步变得更精确的提问。

\hypertarget{ux7b2cux4e09ux5341ux7ae0-ux771fux7406ux4e0eux5171ux540cux6f14ux5316}{%
\chapter{第三十章　真理与共同演化}\label{ux7b2cux4e09ux5341ux7ae0-ux771fux7406ux4e0eux5171ux540cux6f14ux5316}}

一个人可以真诚相信错误的事，一个群体可以一致接受错误的事，一个有利于短期生存的信念也可以与事实不符。由此可见，真理不能只由确信、共识或适应成功定义。本书采取一种开放的现实主义：关于世界的主张之所以为真，取决于世界实际怎样，而我们的认识则通过证据、推理与修正逐步接近它。

这个立场不要求我们已经拥有不受视角限制的全部知识。任何判断都使用某种语言和方法，观察也由仪器、问题与背景认识组织。但是，方法的介入不意味着结果只由我们决定。桥会不会承重，某种响应是否重复出现，一个反例是否满足前提，都可能迫使我们改变期待。世界通过这些不能任意协商的后果进入认识。

真理的客观性与认识的历史性因此可以并存。我们在历史中形成概念，发现遗漏，改变测量和证明方式；在这个过程中，有些旧结论被推翻，有些被限定，有些则因更广的检验而得到加强。历史变化本身不能决定哪项结论正确，但它可以形成更好的求真条件。

哲学与数学的共同演化，在这里具有清楚位置。哲学提出概念关系和普适企图，数学辨明其前提，给出正向结果和反例；模型与实验进一步检查这些前提怎样进入现实。随后，哲学重新理解原来的概念，数学重新处理由新问题引出的结构。两部著作之所以分开，是为了让每一种论证承担自己的责任；它们之所以相互引用，是因为同一问题会穿过这些责任的边界。\footnote{核心第 31 章及附录 C、E；三份科学接口资料；数学著作第 34、38 章。关于客观真理与历史认识的关系，本章采用并辩护开放现实主义立场。}

形式证明的意义也应在这个过程中理解。它使一段推导能够被逐步检查，使``已经给定''与``确实得到''之间的区别更难被修辞遮蔽。它不能替研究者选出所有值得关切的问题，也不能仅凭编译成功确定外部世界符合模型。但当一个看似自然的结论被证明需要额外前提，或者一个小反例揭示了宏大叙述的缺口，哲学便获得了具体的修订理由。

同样，哲学不只是等待技术结果后写一段评论。它可以追问指标是否测量了原来所关切的东西，客体划分是否隐藏了受影响者，目标函数是否偷渡了价值决定，以及某种成功为什么值得称作进步。这些问题能够改变后续模型的设计。一个过于狭窄的问题即使被完美解决，也未必满足原来的认识需要。

共同求真需要的组织，与本书讨论的其他组织一样，具有真实的维持条件。资料要能够检查，批评者要能够发言，错误要能够公开承认，功劳与责任要适当记录。独立复核的价值，不来自另一位研究者必然正确，而来自不同位置能够发现不同遗漏。诚实地保留失败，使知识拥有比成功叙事更完整的记忆。

进步也可以在尚未得到最终答案时发生。把一个模糊命题拆成几个有区分力的问题，证明其中一部分，找到另一部分不能成立的反例，再明确还需怎样的世界证据，这已经实质地改变了我们的处境。知识增长不只表现为肯定句增加，也表现为主张与根据之间的距离缩短。

本书的世界观于是回到开端：客体在变化中形成，通过关系维持，借助记忆和行动参与未来；人的生活在这些自然条件中发展出理由、承诺与共同责任。统一的愿望使我们愿意寻找连接，现实的抵抗使连接必须被认真证明和解释。我们能够维护的，不是一个永远不必修订的思想终点，而是一种让更好的理由、更准确的证据和更公正的共同生活继续出现的能力。

\hypertarget{ux9644ux5f55-a-ux6982ux5ff5ux7684ux4f7fux7528}{%
\chapter{附录 A　概念的使用}\label{ux9644ux5f55-a-ux6982ux5ff5ux7684ux4f7fux7528}}

下表用于澄清本书的术语，不把概念解释当作对所有现实系统的存在证明。

\begin{longtable}[]{@{}
  >{\raggedright\arraybackslash}p{(\columnwidth - 2\tabcolsep) * \real{0.5000}}
  >{\raggedright\arraybackslash}p{(\columnwidth - 2\tabcolsep) * \real{0.5000}}@{}}
\toprule\noalign{}
\begin{minipage}[b]{\linewidth}\raggedright
术语
\end{minipage} & \begin{minipage}[b]{\linewidth}\raggedright
本书采用的含义
\end{minipage} \\
\midrule\noalign{}
\endhead
\bottomrule\noalign{}
\endlastfoot
客体 & 真实过程中，在明确的观察、作用与时间范围内保持可追踪差别的相对稳定组织；其具体解释还要说明实际实现。 \\
过程 & 状态与关系怎样随时间形成、延续和改变；持续客体通过过程实现。 \\
关系性实在论 & 描述随问题选择，成立却受世界的真实联系约束；视角的参与不等于任意创造事实。 \\
边界 & 区分内部与外部作用的接口，可能依预测、干预、资源与身份问题采取不同划分。 \\
Ω & 追问某种组织得以保持所需的构成依赖、冗余与失效方式；几何循环只是可能相关的一类结构。 \\
历史身份 & 由有方向的形成与延续关系支持的同一性，区别于某项任务中的预测等价和功能相似。 \\
开放 & 客体通过交换和外部条件参与持续；开放的作用需按接口、内容、时间与依赖说明。 \\
修复 & 在保存指定连续关系的条件下恢复能力；形成新客体、复制和繁殖需要另外说明。 \\
记忆 & 把历史中的差别保存为后来预测与行动可用的区别；其充分性相对于具体问题。 \\
可行性 & 存在实际能够实施、并对相关后续分支作出回应的完整策略；约束的内容和时限必须给定。 \\
目的 & 使未来结果以成功、失败或偏好差别参与当前行动组织的关系；其功能、历史、意向及共同承诺形式需要分别解释。 \\
Π／Φ & 对实现资源与信息适切性的双重关切；成为定量理论时，需要独立的定义、单位、机制与校准。 \\
评价的表示自由 & 不同数值分解或表示可能保留同样的总评价或选择关系；额外锚定决定能够识别到什么程度。 \\
长期组织 & 不变、循环、占据比例、条件存活及其他指定时间行为；各自需要相应存在和稳定条件。 \\
能动性 & 通过认识、可实行的选择、反馈以及对程序和目标的适当修订，真实地改变未来的能力。 \\
意义 & 主体把活动理解为参与值得关切的事物，并使这种参与进入生活历史的关系；经验投入、现实联系与规范评价各有作用。 \\
自由 & 有理解、回应理由与实际选择机会的有条件能力；与支撑或破坏这些能力的社会安排共同评价。 \\
统一 & 使多种解释之间形成有根据的联系与相容关系；普适单一理论是更强的研究企图。 \\
真理 & 主张与世界实际情况的符合；认识通过历史中的证据、推理与共同修正接近它。 \\
\end{longtable}

\hypertarget{ux9644ux5f55-b-ux4e0eux539fux59cbux8bbaux6587ux7684ux5bf9ux5e94}{%
\chapter{附录 B　与原始论文的对应}\label{ux9644ux5f55-b-ux4e0eux539fux59cbux8bbaux6587ux7684ux5bf9ux5e94}}

原始论文是本书的问题来源。本书按论证关系重新分篇，保留全部主要主题；保留一个问题，不意味着保留原稿关于该问题的全部强断言。页码指 \href{../UEOT_v3_top.pdf}{原始 43 页 PDF} 的实际页序，跨页主题因而会有重叠。

\begin{longtable}[]{@{}
  >{\raggedright\arraybackslash}p{(\columnwidth - 4\tabcolsep) * \real{0.3333}}
  >{\raggedright\arraybackslash}p{(\columnwidth - 4\tabcolsep) * \real{0.3333}}
  >{\raggedright\arraybackslash}p{(\columnwidth - 4\tabcolsep) * \real{0.3333}}@{}}
\toprule\noalign{}
\begin{minipage}[b]{\linewidth}\raggedright
原始论文的位置与主题
\end{minipage} & \begin{minipage}[b]{\linewidth}\raggedright
本书的承接位置
\end{minipage} & \begin{minipage}[b]{\linewidth}\raggedright
主要修订方式
\end{minipage} \\
\midrule\noalign{}
\endhead
\bottomrule\noalign{}
\endlastfoot
第 1---3 页：物质、心灵、意义之间的分裂；系统思想、控制与信息的共同问题 & 前言，第 1、30 章 & 从统一的动机出发，论证解释相容的要求；共同世界与单一终极理论分开。 \\
第 4---5 页：拓扑存在、Ω、客体类型、开放与生命周期 & 第 2---7、20 章 & 以持续组织、构成依赖和失效方式展开；不要求所有客体都出现同一种几何回路或生命周期。 \\
第 6---7 页：Π／Φ、GOD／GOA、耦合与场 & 第 3、11---14 章 & 保留资源与认识的双重关切，补充目的来源、计量与长期行为的论证；具体作用由实际耦合说明。 \\
第 8 页：同层合作竞争、父子层级 & 第 16、18、26 章 & 区分联合活动、父对象形成和公共正当性；整体持续不自动构成成员服从的理由。 \\
第 9---11 页：结构画像、局部偏好指标、五步诊断 & 第 17、28 章 & 把诊断写成问题、接口、校准与可失败解释；比值和方向指标需要独立有效性。 \\
第 12---15 页：生物演化、开拓与优化、谱系延续 & 第 19---20 章 & 分开修复、个体、遗传与选择单位；创新可以改变可能空间，演化方向需具体机制。 \\
第 15---18 页：自我、意识、注意、意志、记忆、同理心与心理失调 & 第 8、15、21---23 章 & 建立有身体、有历史的主体论；功能组织与经验问题相接，具体临床或神经断言留给独立证据。 \\
第 18---24 页：市场、财富集中、泡沫、网络风险、快慢层级 & 第 18、26 章 & 以制度机制解释局部与整体后果；效率、公正、稳定与预测能力分别评价。 \\
第 24---28 页：主体---世界复合、意义、价值、伦理与自由 & 第 23---27 章 & 以承诺、关怀、公共理由和实际机会给出独立规范论证，保留有理由的异议与个人不可替代性。 \\
第 28---29 页：核心概念、数学与测量的开放问题 & 第 10、13、28、30 章 & 把概念锚定、模型适用、有限证据和纠错组织贯穿方法论。 \\
第 30---33 页：量子与粒子拓扑猜想 & 第 29 章 & 保留统一研究的方向，明确从结构类比到物理理论还需可观察量、动力学和数值对应。 \\
第 34---37 页：信息几何、引力、奇点、宇宙学项、黑洞信息与尺度统一 & 第 17、29 章 & 作为原稿的基础物理远景加以讨论；尺度映射、时间过程和相容动力学分别承担解释责任。 \\
第 37---39 页：拓扑数据、粒子共振及宇宙观测提案 & 第 4、28---29 章 & 实验需要区分竞争解释并先确定具体预测；结构迹象不能单独确认整套世界观。 \\
第 40---41 页：统一世界观、目的与意义、终极作用原则 & 前言，第 24、29---30 章 & 保留开放的自然主义企图；意义有现实承受者，终极统一需要逐项取得证据。 \\
第 41---43 页：参考文献、致谢、作者贡献 & 本书署名说明及附录 D & 保留 Degui Qian 为原论文作者。原稿的参考文献在原 PDF 中完整可查，不把未独立复核的引文改作本书的经验依据。 \\
\end{longtable}

\hypertarget{ux9644ux5f55-c-ux4e0eux72ecux7acbux6570ux5b66ux8457ux4f5cux7684ux63a5ux53e3}{%
\chapter{附录 C　与独立数学著作的接口}\label{ux9644ux5f55-c-ux4e0eux72ecux7acbux6570ux5b66ux8457ux4f5cux7684ux63a5ux53e3}}

配套数学著作为 \href{../mathematics/UEOT_MATHEMATICS.md}{《UEOT 数学著作》}。它保留原核心第 0---31 章，并增补层级、校准、客体与目的接口、同源控制、实际占据、有限支持可行性及选后联合风险。读者可以独立阅读任一部书。下表只定位技术背景，不以技术结论代替哲学理由。

\begin{longtable}[]{@{}
  >{\raggedright\arraybackslash}p{(\columnwidth - 4\tabcolsep) * \real{0.3333}}
  >{\raggedright\arraybackslash}p{(\columnwidth - 4\tabcolsep) * \real{0.3333}}
  >{\raggedright\arraybackslash}p{(\columnwidth - 4\tabcolsep) * \real{0.3333}}@{}}
\toprule\noalign{}
\begin{minipage}[b]{\linewidth}\raggedright
哲学问题
\end{minipage} & \begin{minipage}[b]{\linewidth}\raggedright
数学著作中的精确工作
\end{minipage} & \begin{minipage}[b]{\linewidth}\raggedright
阅读入口
\end{minipage} \\
\midrule\noalign{}
\endhead
\bottomrule\noalign{}
\endlastfoot
观察、干预与客体划分怎样相接 & 协议、规范预测状态、相对于声明响应的最小性 & \href{../mathematics/UEOT_MATHEMATICS.md\#ch-2}{第 2 章}、\href{../mathematics/UEOT_MATHEMATICS.md\#ch-4}{第 4 章} \\
组织完整性如何从毁损方式辨明 & 对象证书、极小失效结构、载体与物理干预的条件 & \href{../mathematics/UEOT_MATHEMATICS.md\#ch-12}{第 12 章}、\href{../mathematics/UEOT_MATHEMATICS.md\#ch-13}{第 13 章} \\
同一性怎样跨变化延续 & 历史运输、形成与绑定、客体角色的分离 & \href{../mathematics/UEOT_MATHEMATICS.md\#ch-11}{第 11 章}、\href{../mathematics/UEOT_MATHEMATICS.md\#ch-34}{第 34 章} \\
多种层级何时可以组合 & 八类层级关系、商降落、误差传输及领域适配 & \href{../mathematics/UEOT_MATHEMATICS.md\#ch-15}{第 15 章}、\href{../mathematics/UEOT_MATHEMATICS.md\#ch-32}{第 32 章} \\
为什么同一关系不足以决定数量尺度 & 缺陷重标度、双边校准、复合预算 & \href{../mathematics/UEOT_MATHEMATICS.md\#ch-33}{第 33 章} \\
目的与客体如何汇合 & 给定评价桥的控制、目的等价、同过程异目标 & \href{../mathematics/UEOT_MATHEMATICS.md\#ch-18}{第 18 章}、\href{../mathematics/UEOT_MATHEMATICS.md\#ch-34}{第 34 章}、\href{../mathematics/UEOT_MATHEMATICS.md\#ch-35}{第 35 章} \\
Π／Φ 的独立意义如何建立 & 规范自由、独立锚定、共同二维假设的必要检验 & \href{../mathematics/UEOT_MATHEMATICS.md\#ch-23}{第 23 章}、\href{../mathematics/UEOT_MATHEMATICS.md\#ch-33}{第 33 章} \\
最优性与长期稳定有哪些不同含义 & 约束策略类、最优动作对应、不变律、占据与条件存活 & \href{../mathematics/UEOT_MATHEMATICS.md\#ch-19}{第 19 章}、\href{../mathematics/UEOT_MATHEMATICS.md\#ch-21}{第 21 章} \\
生存与修复怎样在同一个世界模型中连接 & 真实源核、无限路径、约束最优与可信修复机制 & \href{../mathematics/UEOT_MATHEMATICS.md\#ch-35}{第 35 章} \\
记忆怎样保留行动所需的可能性 & 完整观察历史、实际后验支持、策略无关支持递推 & \href{../mathematics/UEOT_MATHEMATICS.md\#ch-37}{第 37 章} \\
局部合格与完整可行有何区别 & 有限时域策略存在的充要条件、可执行关系检查、可达死路 & \href{../mathematics/UEOT_MATHEMATICS.md\#ch-37}{第 37 章} \\
学习怎样接到真实部署风险 & 同策略实际占据、注册校准、Good 上的义务及真实联合律 & \href{../mathematics/UEOT_MATHEMATICS.md\#ch-36}{第 36 章}、\href{../mathematics/UEOT_MATHEMATICS.md\#ch-38}{第 38 章} \\
遗传与谱系为什么不同于同一个体 & 繁殖机制、同源正算子桥、个体与后代关系 & \href{../mathematics/UEOT_MATHEMATICS.md\#ch-25}{第 25 章}、\href{../mathematics/UEOT_MATHEMATICS.md\#ch-26}{第 26 章}、\href{../mathematics/UEOT_MATHEMATICS.md\#ch-34}{第 34 章} \\
认识与价值何处需要进一步说明 & 反身学习、信息价值、意义的功能／意向／规范接口 & \href{../mathematics/UEOT_MATHEMATICS.md\#ch-27}{第 27 章} \\
\end{longtable}

这些数学结果具有各自的空间、概率、时域、接口与可测性条件。有限模型中的策略存在，不等于所有现实系统都有可行策略；一种合法恢复不等于全部程序都能从局部物理信息中产生；相同过程中的目的差异与主观经验、公共正当性的问题，也需要本书的进一步概念和规范论证。

本版第十章引用的学习例子在数学书中完整给出。它由两个注册源行、每行 256 次指定独立采样构成；数据决定的两种动作分别以 \(7/8\) 和 \(1/8\) 被选择，对应实际失败概率 \(1/8\) 与 \(1/4\)。真实训练---选择---部署联合失败概率为 \(9/64\)，通用校准上界为 \(17/32<1\)。前者是该模型的精确值，后者是所用一般保证，二者没有被混称为尖锐同一界。全为真的训练记录具有正概率，并使经验预算达到 1，因此``每个训练结果都满足同一较小预算''在这个实例中确实不成立。这个例子是明确生成模型中的数学实例，尚不是外部现实系统的数据检验。

一般的支持可行性到联合风险桥，则把由初始支持可行性构造的同一策略接入实际占据与联合律，从结论中移除了另行给定的实际未认证占据上界。上述数值例子直接证明其所选策略的零未认证占据；它与一般桥共享接口，但没有被说成正在使用同一个最优学习算法。零未认证占据只消去覆盖罚项，仍允许正的参考失败与模型误差。

数学书新增第 39 章把同一故障源中的程序解释、实际记录、付费资源、子代修复和目的排序接到原消费者。供给仍是设计条件，费用只核算受控写入；实际无限解释器的记录律及可积收益与同一策略连接，资源供给仍不是它内部产生的。新的预测来源从真实过去和保留潜状态推出有限未来实验的联合律、模零最小充分性及响应更新，并保持同一完整历史控制程序的可数观察记录；不同程序记忆不能因此合并。时戳出生记录给同一保留亲本的不同出生事件以可推导身份。连续来源又从有限 Gaussian 律构造一个全时间概率律及其实际连续修改，接入原 OU 二次恢复；它没有以所需 Brownian 证书为输入，也没有借此完成一般随机积分、Itô 或 Girsanov。声明协议类中，任意真实历史过去与 fresh adaptive 未来的最小性已从同一源导出；未覆盖的协议与源类及其他全局义务仍待完成。这些局部来源结论不表示整个 Core 理论已完成形式化，更不表示现实经验已经确认。

这条来源链也开始落实到随机作用本身。数学书从同一噪声源的整个自然过去推导条件矩，再证明同一时间输入的不同分割给出同一实际输出，最后从真正可预测 σ 代数的稠密性构造完整 predictable L² 上的终点积分。这没有把一般 Itô、Girsanov 或所有过程级结论当成已经完成。有限目的资料的正消费者则保留探针后的真实父对象，以同一个返回程序生成完整无限记录；证书错误、执行故障、训练费用与未来评价仍分别承担含义。

同源随机作用已有进一步过程来源：实际截断和条件投影导出原自然过滤上的连续积分；同一输入的真实时间能量与该积分平方之差，又由过去事件上的实际二阶能量和 Fubini 推导为鞅。原 Brownian 分割平方和的极限也有实际网格误差界。它们没有把任意完整积分的分割二次变差、一般 Itô、Girsanov、受控 SDE 或全部维度／过滤来源当成已完成。串行人口中的实际父身份与经验子代比例则接到原谱系与 Price 消费者；混合程序群体不自动共用一个未核对的维持证书。整个理论的来源义务仍继续。

数学著作第 39.22---39.25 节给同一完整输入的实际分割与加权表、完整训练账本，以及混合程序子代与全部执行／读出结果的共同风险。第 39.26---39.29 节进一步从这些实际随机和导出有明确导数界的同输入标量 Itô／Dynkin，从真实联合过去推导完整向量可预测域及均方连续过程，并构造原 OU 条件核和常控制的实际换测度／KL 消费者。均方连续尚不表示样本连续；这些来源不完成一般受控 SDE 或全部可预测控制。数学闭环与现实经验支持继续分列。

同一数值校准也开始承担实际选择。在二维未知正有理对角增益、两个互异且位于 \([-1,1]\) 的有理潜在响应，以及明确的独立读出权限下，原选择器只读取实际有理观测，使用同一个计算出的逆和动作分数。源推导的保证是整个实验中``接受且错误''的联合概率界，不能解释为每一次接受必定正确。一个合法的完整字确实以正概率产生接受错误；它不是被删除的异常。登记的128条符号及合成观测是形式源字的逐项重建，未被称作新增经验采样，也未改变原有真实批次。把真实未知排序留在证明条件而不交给选择器，保留误差与失败，才使证书说明同一个实际决定的作用范围。

\hypertarget{ux9644ux5f55-d-ux6765ux6e90ux8bbaux8bc1ux5c42ux6b21ux4e0eux5f15ux6ce8}{%
\chapter{附录 D　来源、论证层次与引注}\label{ux9644ux5f55-d-ux6765ux6e90ux8bbaux8bc1ux5c42ux6b21ux4e0eux5f15ux6ce8}}

\hypertarget{d.1-ux6838ux5fc3ux6765ux6e90}{%
\section{D.1 核心来源}\label{d.1-ux6838ux5fc3ux6765ux6e90}}

\textbf{原著。} Degui Qian, \emph{A Universal Evolutionary Object Theory (UEOT): A Framework and Research Program for a Unified Science}，43 页。公开历史文本为仓库根目录的 \href{../UEOT_v3_top.pdf}{UEOT\_v3\_top.pdf}，原始 LaTeX 主文保存在 \href{source/original/main.tex}{main.tex}。本版同时核对了这两种文本。原著的思想来源、参考文献、致谢与作者贡献说明继续以其自身文本为准。

\textbf{核心。} \emph{\nolinkurl{UEOT_Core_Mathematics_v3.0_Complete.md}}，原第 0---31 章及附录。它把原始问题重构为协议相对的预测、实现、动态、身份、评价与组合理论。本版从其中吸收概念约束，具体数学由配套著作保留并增修。

\textbf{科学接口资料。} \nolinkurl{CORE_V3_OPEN_PORT_MATRIX.md}、\nolinkurl{CORE_V3_POST_CORE_SCIENTIFIC_CLOSURE.md}、\nolinkurl{THEORY_COMPLETION_ROADMAP_V2.md}，均为 2026 年 10 月 7 日版本。这些资料帮助区分条件定理、模型实现、有限认证和现实支持；它们的历史任务状态不被用于替代后续实际结果。

\textbf{后续形式化与数学著作。} \href{../formalization/ueot-core/}{UEOT Core} 中的层级、缺陷校准、客体性、目的表示、同源控制、修复、统计部署及有限支持可行性结果，为附录 C 的对应提供技术来源。最新陈述、证明和适用条件在独立\href{../mathematics/UEOT_MATHEMATICS.md}{数学著作}中展开，哲学主文采用自然语言讨论其含义。

\hypertarget{d.2-ux672cux4e66ux600eux6837ux4f7fux7528ux8fd9ux4e9bux6765ux6e90}{%
\section{D.2 本书怎样使用这些来源}\label{d.2-ux672cux4e66ux600eux6837ux4f7fux7528ux8fd9ux4e9bux6765ux6e90}}

本书的客体实在论、自然主义、意义观与规范立场是需要哲学辩护的主张。文中的桥梁、走廊、制度分配等例子是为了展开概念而构造的思想例，不冒充经验调查。涉及精确概率的有限学习例子，有另书的完整数学模型支撑。

原文所列神经科学、金融、粒子与宇宙学研究，保留为其原有文献背景；本书没有把其中尚未独立复核的具体经验断言作为新论证的已确认前提。现实中的普适双驱动、意识机制、社会预测和基础物理统一，仍需相应领域的独立检验。这样的区分使哲学能够继续提出广阔问题，同时让读者知道每种主张由什么支持。

\backmatter
\end{document}
